\RequirePackage{plautopatch}
\documentclass[platex,dvipdfmx,a4paper]{jsarticle}

\usepackage{amsmath}
\usepackage{amsfonts}
\usepackage{amssymb}

\usepackage[dvipdfmx]{graphicx}
\usepackage{float}

\usepackage{tikz}
\usepackage{caption}

\usepackage[%
  dvipdfmx,%
  setpagesize=false,%
  bookmarks=true,%
  bookmarksdepth=tocdepth,%
  bookmarksnumbered=true,%
  colorlinks=false,%
  pdftitle={},%
  pdfsubject={},%
  pdfauthor={},%
  pdfkeywords={}%
]{hyperref}
\usepackage{pxjahyper}

\usepackage{framed}
\setlength{\OuterFrameSep}{0pt}

\usepackage{physics2}
\usephysicsmodule{ab}
\usephysicsmodule{ab.braket}

\newcommand{\hb}[2]{\frac{\partial #1}{\partial #2}}
\newcommand{\bbn}[2]{\frac{d #1}{d #2}}
\newcommand{\nhb}[3]{\frac{\partial^{#1} #2}{\partial #3^{#1}}}
\newcommand{\nbbn}[3]{\frac{d^{#1} #2}{d #3^{#1}}}
\newcommand{\vectxy}[2]{\begin{bmatrix}#1\#2\end{bmatrix}}
\newcommand{\vectxyz}[3]{\begin{bmatrix}#1\#2\#3\end{bmatrix}}
\newcommand{\vect}[1]{\boldsymbol{\mathrm{#1}}}

\newcommand{\ssum}[3]{\sum_{#1=#2}^#3}

\newcommand{\veff}{V_{\mathrm{eff}}(r)}

\hypersetup{colorlinks=true}
\hypersetup{linkcolor=red}

\title{熱力学演習解答解説}
\author{25B00729　嶋田 悠音}

\begin{document}
\maketitle
\begingroup
  \hypersetup{hidelinks}
  \tableofcontents
\endgroup
\clearpage

\section{1-1}

状態量
\begin{align*}
L=(2x+1)^2(3x-3y+1)
\end{align*}
を考える。

\subsection{(i)}

まず $x$ で偏微分する。

\begin{align*}
\left(\frac{\partial L}{\partial x}\right)_y
 &=4(2x+1)(3x-3y+1)
 +(2x+1)^2\cdot3 \\
 &= (2x+1)\left[12x-12y+4+6x+3\right] \\
 &= (2x+1)(18x-12y+7)。
\end{align*}

次に $y$ で偏微分する：

\begin{align*}
\left(\frac{\partial L}{\partial y}\right)_x
=(2x+1)^2(-3)
=-3(2x+1)^2。
\end{align*}

\subsection{(ii)}

変数
\begin{align*}
z=x-y
\end{align*}
を導入すると、

\begin{align*}
L=(2x+1)^2(3z+1)
\end{align*}
である。

まず $x$ に関する偏微分は

\begin{align*}
\left(\frac{\partial L}{\partial x}\right)_z
 &=4(2x+1)(3z+1)。
\end{align*}

また、$z$ に関する偏微分は

\begin{align*}
\left(\frac{\partial L}{\partial z}\right)_x
 &=3(2x+1)^2。
\end{align*}

\subsection{(iii)}

$x,z$ を独立変数とすると

\begin{align*}
dL
=
\left(\frac{\partial L}{\partial x}\right)_z dx
+
\left(\frac{\partial L}{\partial z}\right)_x dz。
\end{align*}

ここで

\begin{align*}
z=z(x,y)
\end{align*}
であるから、

\begin{align*}
dz
=
\left(\frac{\partial z}{\partial x}\right)_y dx
+
\left(\frac{\partial z}{\partial y}\right)_x dy。
\end{align*}

これを代入すると、

\begin{align*}
dL
&=
\left(\frac{\partial L}{\partial x}\right)_z dx
+
\left(\frac{\partial L}{\partial z}\right)_x
\left[
\left(\frac{\partial z}{\partial x}\right)_y dx
+
\left(\frac{\partial z}{\partial y}\right)_x dy
\right] \\
&=
\left[
\left(\frac{\partial L}{\partial x}\right)_z
+
\left(\frac{\partial L}{\partial z}\right)_x
\left(\frac{\partial z}{\partial x}\right)_y
\right]dx
+
\left(\frac{\partial L}{\partial z}\right)_x
\left(\frac{\partial z}{\partial y}\right)_x dy。
\end{align*}

一方、

\begin{align*}
dL
=
\left(\frac{\partial L}{\partial x}\right)_y dx
+
\left(\frac{\partial L}{\partial y}\right)_x dy
\end{align*}

でもあるので、$dx$ の係数を比較すれば

\begin{align*}
\left(\frac{\partial L}{\partial x}\right)_y
=
\left(\frac{\partial L}{\partial x}\right)_z
+
\left(\frac{\partial L}{\partial z}\right)_x
\left(\frac{\partial z}{\partial x}\right)_y
\end{align*}

が得られる。

\subsection{(iv)}

\begin{align*}
z=x-y
\end{align*}

より

\begin{align*}
\left(\frac{\partial z}{\partial x}\right)_y
=1。
\end{align*}

したがって右辺は

\begin{align*}
\left(\frac{\partial L}{\partial x}\right)_z
+
\left(\frac{\partial L}{\partial z}\right)_x
&=
4(2x+1)(3z+1)
+
3(2x+1)^2 \\
&=
(2x+1)\left[12z+4+6x+3\right]。
\end{align*}

ここで
\begin{align*}
z=x-y
\end{align*}
を代入すると、

\begin{align*}
(2x+1)
\left[
18x-12y+7
\right]
\end{align*}

となり、

\begin{align*}
\left(\frac{\partial L}{\partial x}\right)_y
=
(2x+1)(18x-12y+7)
\end{align*}

と一致する。

よって式(4)が実際に成立することが確認できた。
\section{1-2}

巨視的状態が二つの独立変数で決まる系を考える。状態量 $Z$ は $X,Y$ の関数
\begin{align*}
Z=Z(X,Y)
\end{align*}
として与えられるものとする。

\subsection{(i)}

まず $Y$ を一定に固定する。このとき $Z$ は $X$ のみの関数とみなせるので、

\begin{align*}
dZ
=
\left(\frac{\partial Z}{\partial X}\right)_Y dX
\end{align*}

が成り立つ。

一方、局所的に逆関数が存在すると仮定すると、

\begin{align*}
dX
=
\left(\frac{\partial X}{\partial Z}\right)_Y dZ
\end{align*}

とも書ける。

後式を前式へ代入すると：

\begin{align*}
dZ
&=
\left(\frac{\partial Z}{\partial X}\right)_Y
\left(\frac{\partial X}{\partial Z}\right)_Y
dZ。
\end{align*}

したがって、

\begin{align*}
\left(\frac{\partial Z}{\partial X}\right)_Y
\left(\frac{\partial X}{\partial Z}\right)_Y
=
1
\end{align*}

であり、

\begin{align*}
\left(\frac{\partial Z}{\partial X}\right)_Y
=
\left[
\left(\frac{\partial X}{\partial Z}\right)_Y
\right]^{-1}
\end{align*}

を得る。

以上より、

\begin{align*}
\boxed{
\left(\frac{\partial Z}{\partial X}\right)_Y
=
\left[
\left(\frac{\partial X}{\partial Z}\right)_Y
\right]^{-1}
}
\end{align*}

が示された。

\subsection{(ii)}

まず、$Z$ を一定に保つと、

\begin{align*}
dZ
=
\left(\frac{\partial Z}{\partial X}\right)_Y dX
+
\left(\frac{\partial Z}{\partial Y}\right)_X dY
=
0
\end{align*}

である。

したがって、

\begin{align*}
\left(\frac{\partial Y}{\partial X}\right)_Z
=
-
\frac{
\left(\frac{\partial Z}{\partial X}\right)_Y
}{
\left(\frac{\partial Z}{\partial Y}\right)_X
}.
\end{align*}

ここで式(5)を用いると、

\begin{align*}
\left(\frac{\partial X}{\partial Z}\right)_Y
=
\frac{1}{
\left(\frac{\partial Z}{\partial X}\right)_Y
}.
\end{align*}

したがって、

\begin{align*}
\left(\frac{\partial Y}{\partial X}\right)_Z
\left(\frac{\partial X}{\partial Z}\right)_Y
\left(\frac{\partial Z}{\partial Y}\right)_X
&=
-
\frac{
\left(\frac{\partial Z}{\partial X}\right)_Y
}{
\left(\frac{\partial Z}{\partial Y}\right)_X
}
\frac{
1
}{
\left(\frac{\partial Z}{\partial X}\right)_Y
}
\left(\frac{\partial Z}{\partial Y}\right)_X \\
&=
-1。
\end{align*}

よって、

\begin{align*}
\left(\frac{\partial Y}{\partial X}\right)_Z
\left(\frac{\partial X}{\partial Z}\right)_Y
\left(\frac{\partial Z}{\partial Y}\right)_X
=
-1
\end{align*}

が示された。

\subsection{(iii)}

物理量 $V$ を

\begin{align*}
V=V(X,Y)
\end{align*}

とする。

$Y$ を一定に固定すると、

\begin{align*}
dV
=
\left(\frac{\partial V}{\partial X}\right)_Y dX
\end{align*}

が成り立つ。

同様に、

\begin{align*}
dZ
=
\left(\frac{\partial Z}{\partial X}\right)_Y dX
\end{align*}

であるから、

\begin{align*}
dX
=
\left(\frac{\partial X}{\partial Z}\right)_Y dZ
\end{align*}

である。

これを $dV$ に代入すると：

\begin{align*}
dV
&=
\left(\frac{\partial V}{\partial X}\right)_Y
\left(\frac{\partial X}{\partial Z}\right)_Y
dZ。
\end{align*}

一方、

\begin{align*}
dV
=
\left(\frac{\partial V}{\partial Z}\right)_Y dZ
\end{align*}

とも表せるので、係数比較から

\begin{align*}
\left(\frac{\partial V}{\partial Z}\right)_Y
=
\left(\frac{\partial V}{\partial X}\right)_Y
\left(\frac{\partial X}{\partial Z}\right)_Y
\end{align*}

を得る。

したがって、

\begin{align*}
\boxed{
\left(\frac{\partial V}{\partial Z}\right)_Y
=
\left(\frac{\partial V}{\partial X}\right)_Y
\left(\frac{\partial X}{\partial Z}\right)_Y
}
\end{align*}

が示された。
\section{1-3}

理想気体について、内部エネルギーを

\begin{align*}
U=nC_VT
\end{align*}

とする。また状態方程式

\begin{align*}
PV=nRT
\end{align*}

が成り立つものとする。

\subsection{(i)}

定積熱容量の定義は

\begin{align*}
C_V
=
\frac{1}{n}
\left(
\frac{\delta Q}{dT}
\right)_V
\end{align*}

である。

熱力学第1法則

\begin{align*}
dU=\delta Q-PdV
\end{align*}

において、定積変化では

\begin{align*}
dV=0
\end{align*}

であるから、

\begin{align*}
\delta Q=dU
\end{align*}

となる。

したがって、

\begin{align*}
\delta Q
=
d(nC_VT)
=
nC_VdT
\end{align*}

より、

\begin{align*}
C_V
=
\frac{1}{n}
\left(
\frac{\delta Q}{dT}
\right)_V
\end{align*}

が確かに成り立つ。

\subsection{(ii)}

定圧熱容量の定義は

\begin{align*}
C_P
=
\frac{1}{n}
\left(
\frac{\delta Q}{dT}
\right)_P
\end{align*}

である。

熱力学第1法則より、

\begin{align*}
\delta Q
=
dU+PdV
\end{align*}

である。

理想気体では

\begin{align*}
dU=nC_VdT
\end{align*}

であり、さらに定圧変化では

\begin{align*}
PV=nRT
\end{align*}

を微分して

\begin{align*}
PdV=nRdT
\end{align*}

を得る。

したがって、

\begin{align*}
\delta Q
&=
nC_VdT+nRdT \\
&=
n(C_V+R)dT。
\end{align*}

よって、

\begin{align*}
C_P
=
C_V+R
\end{align*}

を得る。

以上より、

\begin{align*}
C_P-C_V=R
\end{align*}

が示された。これは Mayer の関係式である。

\subsection{(iii)}

比熱比を

\begin{align*}
\gamma
=
\frac{C_P}{C_V}
\end{align*}

と定義する。

Mayer の関係式

\begin{align*}
C_P-C_V=R
\end{align*}

から、

\begin{align*}
(\gamma-1)C_V
=
R
\end{align*}

である。

したがって、

\begin{align*}
C_V
=
\frac{R}{\gamma-1}
\end{align*}

を得る。

さらに、

\begin{align*}
C_P
=
C_V+R
=
\frac{R}{\gamma-1}+R
=
\frac{\gamma R}{\gamma-1}
\end{align*}

となる。

よって、

\begin{align*}
C_V
=
\frac{R}{\gamma-1},
\qquad
C_P
=
\frac{\gamma R}{\gamma-1}
\end{align*}

である。

\subsection{(iv)}

単原子理想気体では自由度は3である。

等分配則より、

\begin{align*}
U
=
\frac{3}{2}nRT
\end{align*}

であるから、

\begin{align*}
C_V
=
\frac{1}{n}
\frac{dU}{dT}
=
\frac{3}{2}R
\end{align*}

を得る。

Mayer の関係式より、

\begin{align*}
C_P
=
C_V+R
=
\frac{5}{2}R
\end{align*}

である。

したがって、

\begin{align*}
\gamma
=
\frac{C_P}{C_V}
=
\frac{5/2}{3/2}
=
\frac{5}{3}
\end{align*}

となる。

以上より、

\begin{align*}
C_V=\frac{3}{2}R,
\qquad
C_P=\frac{5}{2}R,
\qquad
\gamma=\frac{5}{3}
\end{align*}

を得る。
\section{1-4}

理想気体の可逆断熱変化について、Poisson の関係式を導く。

\subsection{(i)}

可逆断熱変化では外界との熱のやり取りがないので、

\begin{align*}
\delta Q=0
\end{align*}

である。

熱力学第1法則

\begin{align*}
dU=\delta Q-PdV
\end{align*}

より、

\begin{align*}
dU=-PdV
\end{align*}

となる。

理想気体では

\begin{align*}
dU=nC_VdT
\end{align*}

であるから、

\begin{align*}
nC_VdT=-PdV
\end{align*}

を得る。

さらに状態方程式

\begin{align*}
P=\frac{nRT}{V}
\end{align*}

を代入すると、

\begin{align*}
nC_VdT
=
-\frac{nRT}{V}dV
\end{align*}

となる。

両辺を $nT$ で割れば、

\begin{align*}
C_V\frac{dT}{T}
=
-R\frac{dV}{V}
\end{align*}

を得る。

ここで

\begin{align*}
C_V=\frac{R}{\gamma-1}
\end{align*}

を用いると、

\begin{align*}
\frac{dT}{T}
=
-(\gamma-1)\frac{dV}{V}
\end{align*}

となる。

これを積分して、

\begin{align*}
\int \frac{dT}{T}
=
-(\gamma-1)
\int \frac{dV}{V}
\end{align*}

より、

\begin{align*}
\ln T
=
-(\gamma-1)\ln V+\mathrm{const.}
\end{align*}

を得る。

したがって、

\begin{align*}
\ln \left(TV^{\gamma-1}\right)
=
\mathrm{const.}
\end{align*}

であるから、

\begin{align*}
TV^{\gamma-1}
=
\mathrm{const.}
\end{align*}

が得られる。

\subsection{(ii)}

状態方程式

\begin{align*}
PV=nRT
\end{align*}

を用いて温度を消去する。

すなわち、

\begin{align*}
T=\frac{PV}{nR}
\end{align*}

を

\begin{align*}
TV^{\gamma-1}
=
\mathrm{const.}
\end{align*}

へ代入すると、

\begin{align*}
\frac{PV}{nR}V^{\gamma-1}
=
\mathrm{const.}
\end{align*}

となる。

定数 $nR$ を右辺へ吸収すれば、

\begin{align*}
PV^\gamma
=
\mathrm{const.}
\end{align*}

を得る。

これが Poisson の関係式である。

\subsection{(iii)}

さらに Poisson の関係式から

\begin{align*}
V
=
\left(
\frac{\mathrm{const.}}{P}
\right)^{1/\gamma}
\end{align*}

である。

これを

\begin{align*}
TV^{\gamma-1}
=
\mathrm{const.}
\end{align*}

へ代入すると、

\begin{align*}
T
\left(
\frac{\mathrm{const.}}{P}
\right)^{(\gamma-1)/\gamma}
=
\mathrm{const.}
\end{align*}

となる。

定数をまとめ直せば、

\begin{align*}
TP^{-(\gamma-1)/\gamma}
=
\mathrm{const.}
\end{align*}

すなわち、

\begin{align*}
T^\gamma P^{1-\gamma}
=
\mathrm{const.}
\end{align*}

を得る。

したがって、

\begin{align*}
PV^\gamma
=
\mathrm{const.},
\qquad
TV^{\gamma-1}
=
\mathrm{const.},
\qquad
T^\gamma P^{1-\gamma}
=
\mathrm{const.}
\end{align*}

が互いに等価な Poisson の関係式である。

\subsection{(iv)}

可逆断熱変化において体積が

\begin{align*}
V_1 \rightarrow V_2
\end{align*}

へ変化するとする。

Poisson の関係式より、

\begin{align*}
P
=
KV^{-\gamma}
\end{align*}

と書ける。ただし $K$ は定数である。

気体が外界へする仕事は

\begin{align*}
W
=
\int_{V_1}^{V_2} P\,dV
\end{align*}

であるから、

\begin{align*}
W
&=
K\int_{V_1}^{V_2}V^{-\gamma}dV \\
&=
\frac{K}{1-\gamma}
\left(
V_2^{\,1-\gamma}
-
V_1^{\,1-\gamma}
\right)。
\end{align*}

ここで

\begin{align*}
K=P_1V_1^\gamma
=P_2V_2^\gamma
\end{align*}

を用いると、

\begin{align*}
KV_1^{1-\gamma}
=
P_1V_1,
\qquad
KV_2^{1-\gamma}
=
P_2V_2
\end{align*}

である。

したがって、

\begin{align*}
W
=
\frac{P_1V_1-P_2V_2}{\gamma-1}
\end{align*}

を得る。

さらに理想気体の状態方程式より、

\begin{align*}
P_1V_1=nRT_1,
\qquad
P_2V_2=nRT_2
\end{align*}

であるから、

\begin{align*}
W
=
\frac{nR(T_1-T_2)}{\gamma-1}
\end{align*}

となる。

ここで

\begin{align*}
C_V=\frac{R}{\gamma-1}
\end{align*}

を用いれば、

\begin{align*}
W
=
nC_V(T_1-T_2)
\end{align*}

を得る。

これは断熱変化では

\begin{align*}
Q=0
\end{align*}

であり、

\begin{align*}
W=-\Delta U
\end{align*}

となることと一致している。

以上より、可逆断熱変化における Poisson の関係式および仕事の表式が示された。
\section{1-5}

断熱材に覆われたピストン中の 1 mol の単原子理想気体を、準静的かつ断熱的に

\begin{align*}
V_1=2.00\,\mathrm{m^3}
\end{align*}

から

\begin{align*}
V_2=1.00\,\mathrm{m^3}
\end{align*}

まで圧縮する。

比熱比は

\begin{align*}
\gamma=\frac53
\end{align*}

であり、初期温度は

\begin{align*}
T_1=300\,\mathrm{K}
\end{align*}

とする。

\subsection{(1)}

準静的断熱変化では Poisson の関係式

\begin{align*}
TV^{\gamma-1}
=
\mathrm{const.}
\end{align*}

が成り立つ。

したがって、

\begin{align*}
T_1V_1^{\gamma-1}
=
T_2V_2^{\gamma-1}
\end{align*}

より、

\begin{align*}
T_2
=
T_1
\left(
\frac{V_1}{V_2}
\right)^{\gamma-1}
\end{align*}

を得る。

ここで、

\begin{align*}
\gamma-1
=
\frac23
\end{align*}

であるから、

\begin{align*}
T_2
&=
300\times
2^{2/3} \\
&=
300\times1.58740 \\
&=
476.220\cdots\,\mathrm{K}。
\end{align*}

よって、有効数字3桁で

\begin{align*}
T_2
=
4.76\times10^2\,\mathrm{K}
\end{align*}

である。

\subsection{(2)}

断熱過程では

\begin{align*}
Q=0
\end{align*}

である。

熱力学第1法則を

\begin{align*}
\Delta U=W+Q
\end{align*}

と書くと、

\begin{align*}
W=\Delta U
\end{align*}

となる。ただし問題文の記号 $W$ は「外界が気体にした仕事」である。

単原子理想気体の内部エネルギーは

\begin{align*}
U=\frac32NRT
\end{align*}

であるから、

\begin{align*}
W
&=
\frac32NR(T_2-T_1)。
\end{align*}

ここで

\begin{align*}
N=1.00\,\mathrm{mol},
\qquad
R=8.314\,\mathrm{J\,mol^{-1}K^{-1}}
\end{align*}

および

\begin{align*}
T_2-T_1
=
476.220-300
=
176.220\,\mathrm{K}
\end{align*}

を代入すると、

\begin{align*}
W
&=
\frac32
\times
8.314
\times
176.220 \\
&=
2.20\times10^3\,\mathrm{J}。
\end{align*}

したがって、ピストンを押すために外界が気体へ行った仕事は

\begin{align*}
W
=
2.20\times10^3\,\mathrm{J}
\end{align*}

である。

以上より、

\begin{align*}
T_2
=
4.76\times10^2\,\mathrm{K},
\qquad
W
=
2.20\times10^3\,\mathrm{J}
\end{align*}

を得る。
\section{1-6}

\subsection{(i)}

華厳の滝の落差を

\begin{align*}
h=97.0\,\mathrm{m}
\end{align*}

とする。

質量 $m$ の水塊が落下するとき、失われる位置エネルギーは

\begin{align*}
\Delta E
=
mgh
\end{align*}

である。

問題文の仮定より、このエネルギーがすべて熱へ変換されるとする。

一方、水の温度が $\Delta T$ 上昇するときに吸収する熱量は

\begin{align*}
Q=mc\Delta T
\end{align*}

である。

エネルギー保存則より、

\begin{align*}
mgh
=
mc\Delta T
\end{align*}

が成り立つので、

\begin{align*}
\Delta T
=
\frac{gh}{c}
\end{align*}

を得る。

ここで、比熱

\begin{align*}
c
=
1.00\,\mathrm{cal\,g^{-1}K^{-1}}
\end{align*}

を SI 単位へ変換する。

熱の仕事当量

\begin{align*}
1\,\mathrm{cal}
=
4.184\,\mathrm{J}
\end{align*}

より、

\begin{align*}
c
=
4.184\times10^3
\,\mathrm{J\,kg^{-1}K^{-1}}
\end{align*}

である。

したがって、

\begin{align*}
\Delta T
&=
\frac{9.81\times97.0}
{4.184\times10^3}
\\
&=
2.27\times10^{-1}\,\mathrm{K}。
\end{align*}

初期温度は

\begin{align*}
T_1=15.0\,^\circ\mathrm{C}
\end{align*}

であるから、

\begin{align*}
T_2
&=
15.0+0.227 \\
&=
15.227\,^\circ\mathrm{C}。
\end{align*}

よって、有効数字3桁で

\begin{align*}
T_2
=
15.2\,^\circ\mathrm{C}
\end{align*}

となる。

\subsection{(ii)}

静止した大気を考える。

高さ方向に厚さ $dh$、断面積 $A$ の微小な空気柱を取る。

この空気柱に働く力のつり合いより、

\begin{align*}
p(h)A-p(h+dh)A
=
\rho gA\,dh
\end{align*}

が成り立つ。

両辺を $A\,dh$ で割り、

\begin{align*}
\frac{p(h+dh)-p(h)}{dh}
=
-\rho g
\end{align*}

と書く。

極限 $dh\to0$ を取れば、

\begin{align*}
\frac{dp}{dh}
=
-\rho g
\end{align*}

を得る。

ここで理想気体の状態方程式

\begin{align*}
pV=nRT
\end{align*}

を用いる。

モル質量を $m$ とすると、

\begin{align*}
\rho
=
\frac{nm}{V}
\end{align*}

である。

また、

\begin{align*}
\frac{n}{V}
=
\frac{p}{RT}
\end{align*}

であるから、

\begin{align*}
\rho
=
\frac{mp}{RT}
\end{align*}

を得る。

したがって、

\begin{align*}
\frac{dp}{dh}
=
-\frac{mpg}{RT}
\end{align*}

となる。

以上より、

\begin{align*}
\frac{dp}{dh}
=
-\frac{mpg}{RT}
\end{align*}

が示された。

\subsection{(iii)}

Poisson の関係式

\begin{align*}
pV^\gamma
=
\mathrm{const.}
\end{align*}

の対数を取る：

\begin{align*}
d(\ln p)
+
\gamma d(\ln V)
=
0。
\end{align*}

したがって、

\begin{align*}
\frac{dp}{p}
+
\gamma\frac{dV}{V}
=
0
\end{align*}

である。

一方、理想気体の状態方程式

\begin{align*}
pV=RT
\end{align*}

を微分すると、

\begin{align*}
\frac{dp}{p}
+
\frac{dV}{V}
=
\frac{dT}{T}
\end{align*}

を得る。

最初の式から

\begin{align*}
\frac{dV}{V}
=
-\frac1\gamma
\frac{dp}{p}
\end{align*}

であるので、

\begin{align*}
\frac{dT}{T}
&=
\left(
1-\frac1\gamma
\right)
\frac{dp}{p}
\\
&=
\frac{\gamma-1}{\gamma}
\frac{dp}{p}
\end{align*}

となる。

したがって、

\begin{align*}
\frac{dT}{dh}
=
\frac{\gamma-1}{\gamma}
\frac{T}{p}
\frac{dp}{dh}
\end{align*}

を得る。

前問の結果

\begin{align*}
\frac{dp}{dh}
=
-\frac{mpg}{RT}
\end{align*}

を代入すると、

\begin{align*}
\frac{dT}{dh}
&=
\frac{\gamma-1}{\gamma}
\frac{T}{p}
\left(
-\frac{mpg}{RT}
\right)
\\
&=
-\frac{mg}{R}
\frac{\gamma-1}{\gamma}。
\end{align*}

よって、

\begin{align*}
\frac{dT}{dh}
=
-\frac{mg}{R}
\frac{\gamma-1}{\gamma}
\end{align*}

が示された。

\subsection{(iv)}

空気について

\begin{align*}
m
=
28.97\times10^{-3}
\,\mathrm{kg\,mol^{-1}},
\qquad
\gamma=\frac75
\end{align*}

を用いる。

前問の結果より、

\begin{align*}
\frac{dT}{dh}
&=
-\frac{mg}{R}
\frac{\gamma-1}{\gamma}
\\
&=
-\frac{
(28.97\times10^{-3})
\times
9.81
}{
8.314
}
\times
\frac{2}{7}
\\
&=
-9.77\times10^{-3}
\,\mathrm{K\,m^{-1}}。
\end{align*}

したがって、

\begin{align*}
\Delta h
=
97.0\,\mathrm{m}
\end{align*}

について

\begin{align*}
\Delta T
&=
\frac{dT}{dh}\Delta h
\\
&=
-(9.77\times10^{-3})
\times97.0
\\
&=
-9.48\times10^{-1}
\,\mathrm{K}。
\end{align*}

よって、

\begin{align*}
\Delta T
=
-0.948\,^\circ\mathrm{C}
\end{align*}

である。

すなわち、華厳の滝程度の高度差では、断熱膨張による気温低下はおよそ

\begin{align*}
0.95\,^\circ\mathrm{C}
\end{align*}

である。

一方、(i) で求めた水温上昇は

\begin{align*}
0.227\,^\circ\mathrm{C}
\end{align*}

であった。

したがって、

\begin{align*}
0.95
>
0.23
\end{align*}

であり、気温の高度依存による効果の方が約4倍大きい。

以上より、華厳の滝では落下による水温上昇よりも、大気の断熱変化による気温低下の方が支配的であることが分かる。
\section{2-1}

熱容量を一般の準静的過程に沿って考える。熱力学第1法則と第1回で導出した定圧熱容量・定積熱容量の関係式を用いて、一般の熱容量を表す。

\subsection{(i)}

準静的過程 $q$ に沿った熱容量を

\begin{align*}
C_q
=
\frac{\delta Q}{dT}
\end{align*}

と定義する。

熱力学第1法則より、

\begin{align*}
\delta Q
=
dU+p\,dV
\end{align*}

である。

内部エネルギーを変数 $T,V$ の関数として表すと、

\begin{align*}
dU
=
\left(\frac{\partial U}{\partial T}\right)_V dT
+
\left(\frac{\partial U}{\partial V}\right)_T dV
\end{align*}

であるから、

\begin{align*}
\delta Q
=
\left(\frac{\partial U}{\partial T}\right)_V dT
+
\left[
\left(\frac{\partial U}{\partial V}\right)_T+p
\right]dV。
\end{align*}

ここで

\begin{align*}
C_V
=
\left(\frac{\partial U}{\partial T}\right)_V
\end{align*}

を用いると、

\begin{align*}
\delta Q
=
C_V dT
+
\left[
\left(\frac{\partial U}{\partial V}\right)_T+p
\right]dV。
\end{align*}

一方、第1回の結果より、

\begin{align*}
C_P
=
C_V+
\left[
\left(\frac{\partial U}{\partial V}\right)_T+p
\right]
\left(\frac{\partial V}{\partial T}\right)_p
\end{align*}

である。

したがって、

\begin{align*}
\left[
\left(\frac{\partial U}{\partial V}\right)_T+p
\right]
=
\frac{C_P-C_V}
{\left(\frac{\partial V}{\partial T}\right)_p}
\end{align*}

を得る。

これを先ほどの式へ代入すると、

\begin{align*}
\delta Q
=
C_VdT
+
(C_P-C_V)
\frac{
\left(\frac{\partial V}{\partial T}\right)_q
}{
\left(\frac{\partial V}{\partial T}\right)_p
}
dT。
\end{align*}

よって、

\begin{align*}
C_q
=
C_V
+
(C_P-C_V)
\frac{
\left(\frac{\partial V}{\partial T}\right)_q
}{
\left(\frac{\partial V}{\partial T}\right)_p
}
\end{align*}

となる。

膨張率

\begin{align*}
\alpha
=
\frac1V
\left(\frac{\partial V}{\partial T}\right)_p
\end{align*}

を用いると、

\begin{align*}
\left(\frac{\partial V}{\partial T}\right)_p
=
\alpha V
\end{align*}

であるから、

\begin{align*}
C_q
=
C_V
+
\frac{C_P-C_V}
{\alpha V}
\left(\frac{\partial V}{\partial T}\right)_q
\end{align*}

を得る。

以上より、

\begin{align*}
C_q
=
C_V
+
\frac{C_P-C_V}
{\alpha V}
\left(\frac{\partial V}{\partial T}\right)_q
\end{align*}

が示された。

\subsection{(ii)}

圧力 $p$ と体積 $V$ を独立変数に取る。

内部エネルギーの全微分は

\begin{align*}
dU
=
\left(\frac{\partial U}{\partial p}\right)_V dp
+
\left(\frac{\partial U}{\partial V}\right)_p dV
\end{align*}

である。

断熱過程では

\begin{align*}
\delta Q=0
\end{align*}

であるから、熱力学第1法則より

\begin{align*}
dU+p\,dV=0
\end{align*}

となる。

したがって、

\begin{align*}
\left(\frac{\partial U}{\partial p}\right)_V dp
+
\left[
\left(\frac{\partial U}{\partial V}\right)_p+p
\right]dV
=
0。
\end{align*}

これを変形すると、

\begin{align*}
\left(\frac{\partial U}{\partial p}\right)_V dp
=
-
\left[
\left(\frac{\partial U}{\partial V}\right)_p+p
\right]dV。
\end{align*}

よって、断熱条件下での微分比は

\begin{align*}
\left(\frac{\partial p}{\partial V}\right)_{\mathrm{adia}}
=
-
\frac{
\left(\frac{\partial U}{\partial V}\right)_p+p
}{
\left(\frac{\partial U}{\partial p}\right)_V
}
\end{align*}

となる。

以上より、

\begin{align*}
\left(\frac{\partial p}{\partial V}\right)_{\mathrm{adia}}
=
-
\frac{
\left(\frac{\partial U}{\partial V}\right)_p+p
}{
\left(\frac{\partial U}{\partial p}\right)_V
}
\end{align*}

が示された。
\section{2-2}

熱容量と体積弾性率の関係を調べる。

\subsection{(i)}

前回の演習で得られた定圧熱容量の式

\begin{align*}
C_P
=
\left(\frac{\partial U}{\partial T}\right)_V
+
\left[
\left(\frac{\partial U}{\partial V}\right)_T+p
\right]
\left(\frac{\partial V}{\partial T}\right)_p
\end{align*}

を用いる。

ここで内部エネルギーを変数 $p,V$ の関数として考える。

全微分は

\begin{align*}
dU
=
\left(\frac{\partial U}{\partial p}\right)_Vdp
+
\left(\frac{\partial U}{\partial V}\right)_pdV
\end{align*}

である。

定圧変化では

\begin{align*}
dp=0
\end{align*}

であるから

\begin{align*}
dU
=
\left(\frac{\partial U}{\partial V}\right)_p dV
\end{align*}

となる。

一方、

\begin{align*}
dU
=
\left(\frac{\partial U}{\partial T}\right)_p dT
\end{align*}

でもあるので、

\begin{align*}
\left(\frac{\partial U}{\partial T}\right)_p
=
\left(\frac{\partial U}{\partial V}\right)_p
\left(\frac{\partial V}{\partial T}\right)_p
\end{align*}

を得る。

したがって、

\begin{align*}
C_P
&=
\left(\frac{\partial U}{\partial V}\right)_p
\left(\frac{\partial V}{\partial T}\right)_p
+
\left[
\left(\frac{\partial U}{\partial V}\right)_T+p
\right]
\left(\frac{\partial V}{\partial T}\right)_p \\
&=
\left[
\left(\frac{\partial U}{\partial V}\right)_p+p
\right]
\left(\frac{\partial V}{\partial T}\right)_p。
\end{align*}

以上より、

\begin{align*}
C_P
=
\left[
\left(\frac{\partial U}{\partial V}\right)_p+p
\right]
\left(\frac{\partial V}{\partial T}\right)_p
\end{align*}

が示された。

\subsection{(ii)}

まず、

\begin{align*}
k_T
=
-
V
\left(\frac{\partial p}{\partial V}\right)_T,
\qquad
k_{\mathrm{adia}}
=
-
V
\left(\frac{\partial p}{\partial V}\right)_{\mathrm{adia}}
\end{align*}

である。

したがって、

\begin{align*}
\frac{k_{\mathrm{adia}}}{k_T}
=
\frac{
\left(\frac{\partial p}{\partial V}\right)_{\mathrm{adia}}
}{
\left(\frac{\partial p}{\partial V}\right)_T
}
\end{align*}

となる。

ここで前問 2-1(ii) より、

\begin{align*}
\left(\frac{\partial p}{\partial V}\right)_{\mathrm{adia}}
=
-
\frac{
\left(\frac{\partial U}{\partial V}\right)_p+p
}{
\left(\frac{\partial U}{\partial p}\right)_V
}
\end{align*}

である。

一方、

\begin{align*}
C_V
=
\left(\frac{\partial U}{\partial T}\right)_V
\end{align*}

であり、

\begin{align*}
dU
=
\left(\frac{\partial U}{\partial p}\right)_Vdp
\end{align*}

から、

\begin{align*}
\left(\frac{\partial U}{\partial p}\right)_V
=
C_V
\left(\frac{\partial T}{\partial p}\right)_V
\end{align*}

を得る。

したがって、

\begin{align*}
\left(\frac{\partial p}{\partial V}\right)_{\mathrm{adia}}
=
-
\frac{
\left(\frac{\partial U}{\partial V}\right)_p+p
}{
C_V
\left(\frac{\partial T}{\partial p}\right)_V
}
\end{align*}

である。

さらに (i) の結果から、

\begin{align*}
\left(\frac{\partial U}{\partial V}\right)_p+p
=
\frac{
C_P
}{
\left(\frac{\partial V}{\partial T}\right)_p
}
\end{align*}

であるので、

\begin{align*}
\left(\frac{\partial p}{\partial V}\right)_{\mathrm{adia}}
=
-
\frac{C_P}
{C_V}
\frac{1}
{
\left(\frac{\partial V}{\partial T}\right)_p
\left(\frac{\partial T}{\partial p}\right)_V
}。
\end{align*}

ここで循環関係

\begin{align*}
\left(\frac{\partial V}{\partial T}\right)_p
\left(\frac{\partial T}{\partial p}\right)_V
\left(\frac{\partial p}{\partial V}\right)_T
=
-1
\end{align*}

を用いると、

\begin{align*}
\frac{1}
{
\left(\frac{\partial V}{\partial T}\right)_p
\left(\frac{\partial T}{\partial p}\right)_V
}
=
-
\left(\frac{\partial p}{\partial V}\right)_T
\end{align*}

である。

したがって、

\begin{align*}
\left(\frac{\partial p}{\partial V}\right)_{\mathrm{adia}}
=
\frac{C_P}{C_V}
\left(\frac{\partial p}{\partial V}\right)_T
\end{align*}

を得る。

よって、

\begin{align*}
\frac{k_{\mathrm{adia}}}{k_T}
&=
\frac{
\left(\frac{\partial p}{\partial V}\right)_{\mathrm{adia}}
}{
\left(\frac{\partial p}{\partial V}\right)_T
}
\\
&=
\frac{C_P}{C_V}。
\end{align*}

以上より、

\begin{align*}
\frac{k_{\mathrm{adia}}}{k_T}
=
\frac{C_P}{C_V}
\end{align*}

が示された。
\section{2-3}

理想気体のエントロピーおよび内部エネルギーの基本式を導く。

\subsection{(i)}

まずエントロピーの微分形を求める。

熱力学の基本関係式

\begin{align*}
dU=TdS-pdV
\end{align*}

より、

\begin{align*}
dS
=
\frac{dU}{T}
+\frac{p}{T}dV
\end{align*}

を得る。

理想気体では内部エネルギーは温度のみの関数であり、

\begin{align*}
dU=C_VdT
\end{align*}

である。

また、状態方程式

\begin{align*}
pV=RNT
\end{align*}

より、

\begin{align*}
\frac{p}{T}
=
\frac{RN}{V}
\end{align*}

である。

したがって、

\begin{align*}
dS
&=
\frac{C_V}{T}dT
+
\frac{RN}{V}dV。
\end{align*}

以上より、

\begin{align*}
dS
=
\frac{C_V}{T}dT
+
\frac{RN}{V}dV
\end{align*}

が示された。

\subsection{(ii)}

まず状態 I$(T_0,V_0)$ から状態 II$(T,V_0)$ まで定積変化させる。

この区間では

\begin{align*}
dV=0
\end{align*}

であるから、

\begin{align*}
dS
=
\frac{C_V}{T}dT
\end{align*}

となる。

したがって、

\begin{align*}
S_{\mathrm{II}}-S_{\mathrm{I}}
&=
\int_{T_0}^{T}
\frac{C_V}{T'}\,dT' \\
&=
C_V
\log\frac{T}{T_0}。
\end{align*}

次に状態 II から状態 III$(T,V)$ まで等温変化させる。

この区間では

\begin{align*}
dT=0
\end{align*}

であるから、

\begin{align*}
dS
=
RN\frac{dV}{V}
\end{align*}

となる。

よって、

\begin{align*}
S_{\mathrm{III}}-S_{\mathrm{II}}
&=
RN
\int_{V_0}^{V}
\frac{dV'}{V'} \\
&=
RN
\log\frac{V}{V_0}。
\end{align*}

以上を加えると、

\begin{align*}
S_{\mathrm{III}}-S_{\mathrm{I}}
&=
(S_{\mathrm{III}}-S_{\mathrm{II}})
+(S_{\mathrm{II}}-S_{\mathrm{I}}) \\
&=
C_V\log\frac{T}{T_0}
+
RN\log\frac{V}{V_0}。
\end{align*}

したがって、

\begin{align*}
S_{\mathrm{III}}-S_{\mathrm{I}}
=
C_V\log\frac{T}{T_0}
+
RN\log\frac{V}{V_0}
\end{align*}

が得られる。

\subsection{(iii)}

今度は

\begin{align*}
\mathrm{I}(T_0,V_0)
\rightarrow
\mathrm{IV}(T_0,V)
\rightarrow
\mathrm{III}(T,V)
\end{align*}

という経路を考える。

まず状態 I から状態 IV までは等温変化である。

したがって、

\begin{align*}
S_{\mathrm{IV}}-S_{\mathrm{I}}
=
RN\log\frac{V}{V_0}。
\end{align*}

次に状態 IV から状態 III までは定積変化である。

よって、

\begin{align*}
S_{\mathrm{III}}-S_{\mathrm{IV}}
=
C_V\log\frac{T}{T_0}。
\end{align*}

以上より、

\begin{align*}
S_{\mathrm{III}}-S_{\mathrm{I}}
&=
(S_{\mathrm{III}}-S_{\mathrm{IV}})
+(S_{\mathrm{IV}}-S_{\mathrm{I}}) \\
&=
C_V\log\frac{T}{T_0}
+
RN\log\frac{V}{V_0}。
\end{align*}

これは小問(ii)で得た結果と完全に一致する。

したがって、エントロピー差は経路によらず始状態と終状態だけで決まることが確認できた。

\subsection{(iv)}

小問(ii)の結果を用いる。

基準状態

\begin{align*}
(S_0,T_0,V_0)
\end{align*}

を導入すると、

\begin{align*}
S-S_0
=
C_V\log\frac{T}{T_0}
+
RN\log\frac{V}{V_0}
\end{align*}

である。

ここで

\begin{align*}
C_V=cRN
\end{align*}

と書くと、

\begin{align*}
S-S_0
=
cRN\log\frac{T}{T_0}
+
RN\log\frac{V}{V_0}。
\end{align*}

両辺を $RN$ で割る：

\begin{align*}
\frac{S-S_0}{RN}
=
c\log\frac{T}{T_0}
+
\log\frac{V}{V_0}。
\end{align*}

したがって、

\begin{align*}
c\log\frac{T}{T_0}
=
\frac{S-S_0}{RN}
-\log\frac{V}{V_0}
\end{align*}

であり、

\begin{align*}
\log\frac{T}{T_0}
=
\frac{S-S_0}{cRN}
-\frac{1}{c}\log\frac{V}{V_0}
\end{align*}

を得る。

指数関数を取ると、

\begin{align*}
T
=
T_0
\exp\left(
\frac{S-S_0}{cRN}
\right)
\left(
\frac{V_0}{V}
\right)^{1/c}。
\end{align*}

一方、理想気体では

\begin{align*}
U=cRNT
\end{align*}

であるから、

\begin{align*}
U
&=
cRNT_0
\exp\left(
\frac{S-S_0}{cRN}
\right)
\left(
\frac{V_0}{V}
\right)^{1/c}。
\end{align*}

ここで

\begin{align*}
U_0=cRNT_0
\end{align*}

を用いると、

\begin{align*}
U(S,V)
=
U_0
\left(
\frac{V_0}{V}
\right)^{1/c}
\exp\left(
\frac{S-S_0}{cRN}
\right)
\end{align*}

を得る。

以上より、

\begin{align*}
U(S,V)
=
U_0
\left(
\frac{V_0}{V}
\right)^{1/c}
\exp\left(
\frac{S-S_0}{cRN}
\right)
\end{align*}

が示された。
\section{2-4}

前問で得られた理想気体の内部エネルギーの基本式

\begin{align*}
U(S,V)
=
U_0
\left(
\frac{V_0}{V}
\right)^{1/c}
\exp\left(
\frac{S-S_0}{cRN}
\right)
\end{align*}

を用いる。

また、

\begin{align*}
p
=
-\left(\frac{\partial U}{\partial V}\right)_S,
\qquad
T
=
\left(\frac{\partial U}{\partial S}\right)_V
\end{align*}

によって圧力と温度を定義する。

\subsection{(i)}

まず温度を求める。

エントロピー $S$ を固定して偏微分すると、

\begin{align*}
T
&=
\left(
\frac{\partial U}{\partial S}
\right)_V \\
&=
U_0
\left(
\frac{V_0}{V}
\right)^{1/c}
\exp\left(
\frac{S-S_0}{cRN}
\right)
\frac{1}{cRN} \\
&=
\frac{U}{cRN}。
\end{align*}

したがって

\begin{align*}
U=cRNT
\end{align*}

を得る。

次に圧力を求める。

体積 $V$ について偏微分すると、

\begin{align*}
\left(
\frac{\partial U}{\partial V}
\right)_S
&=
-\frac1c
U_0
\left(
\frac{V_0}{V}
\right)^{1/c}
\exp\left(
\frac{S-S_0}{cRN}
\right)
\frac1V \\
&=
-\frac{U}{cV}
\end{align*}

となるので、

\begin{align*}
p
=
-\left(
\frac{\partial U}{\partial V}
\right)_S
=
\frac{U}{cV}
\end{align*}

である。

ここで

\begin{align*}
U=cRNT
\end{align*}

を代入すると、

\begin{align*}
p
=
\frac{cRNT}{cV}
=
\frac{RNT}{V}
\end{align*}

となる。

よって、

\begin{align*}
pV=RNT
\end{align*}

を得る。

以上より、理想気体の状態方程式が導かれた。

\subsection{(ii)}

小問(i)で

\begin{align*}
T
=
\frac{U}{cRN}
\end{align*}

を得た。

これを変形すれば、

\begin{align*}
U=cRNT
\end{align*}

となる。

したがって、内部エネルギーは温度のみに比例し、

\begin{align*}
U=cRNT
\end{align*}

と表される。

\subsection{(iii)}

第1回で導いた結果より、

\begin{align*}
C_V=cRN
\end{align*}

である。

また、Mayer の関係式

\begin{align*}
C_P-C_V=RN
\end{align*}

が成り立つので、

\begin{align*}
C_P
=
(c+1)RN
\end{align*}

となる。

したがって熱容量比

\begin{align*}
\gamma
=
\frac{C_P}{C_V}
=
\frac{(c+1)RN}{cRN}
=
\frac{c+1}{c}
\end{align*}

である。

これを $c$ について解くと、

\begin{align*}
\gamma c
=
c+1
\end{align*}

より、

\begin{align*}
c(\gamma-1)
=
1
\end{align*}

となる。

したがって、

\begin{align*}
c
=
\frac{1}{\gamma-1}
\end{align*}

を得る。

以上より、

\begin{align*}
U
=
\frac{RNT}{\gamma-1}
\end{align*}

であり、

\begin{align*}
c
=
\frac{1}{\gamma-1}
\end{align*}

であることが示された。
\section{2-5}

van der Waals 気体の内部エネルギーが

\begin{align*}
U(S,V)
=
U_0
\left(
\frac{V_0}{V-bN}
\right)^{1/c}
\exp\left(
\frac{S-S_0}{cRN}
\right)
-\frac{aN^2}{V}
\end{align*}

で与えられるとする。

\subsection{(i)}

まず温度を求める。

\begin{align*}
T
=
\left(
\frac{\partial U}{\partial S}
\right)_V
\end{align*}

より、

\begin{align*}
T
&=
U_0
\left(
\frac{V_0}{V-bN}
\right)^{1/c}
\exp\left(
\frac{S-S_0}{cRN}
\right)
\frac{1}{cRN}。
\end{align*}

したがって、

\begin{align*}
U+\frac{aN^2}{V}
=
cRNT
\end{align*}

を得る。

次に圧力を求める。

\begin{align*}
p
=
-\left(
\frac{\partial U}{\partial V}
\right)_S
\end{align*}

である。

第1項を微分すると、

\begin{align*}
\frac{\partial}{\partial V}
\left[
U_0
\left(
\frac{V_0}{V-bN}
\right)^{1/c}
\exp\left(
\frac{S-S_0}{cRN}
\right)
\right]
=
-\frac{1}{c(V-bN)}
\left(
U+\frac{aN^2}{V}
\right)
\end{align*}

となる。

また、

\begin{align*}
\frac{\partial}{\partial V}
\left(
-\frac{aN^2}{V}
\right)
=
\frac{aN^2}{V^2}
\end{align*}

である。

したがって、

\begin{align*}
p
&=
\frac{1}{c(V-bN)}
\left(
U+\frac{aN^2}{V}
\right)
-\frac{aN^2}{V^2} \\
&=
\frac{RNT}{V-bN}
-\frac{aN^2}{V^2}。
\end{align*}

整理すると、

\begin{align*}
p+\frac{aN^2}{V^2}
=
\frac{RNT}{V-bN}
\end{align*}

であるから、

\begin{align*}
\left(
p+\frac{aN^2}{V^2}
\right)
(V-bN)
=
NRT
\end{align*}

を得る。

以上より、

\begin{align*}
\left(
p+\frac{aN^2}{V^2}
\right)
(V-bN)
=
NRT
\end{align*}

が導かれた。

\subsection{(ii)}

小問(i)より、

\begin{align*}
U+\frac{aN^2}{V}
=
cRNT
\end{align*}

である。

したがって内部エネルギーは

\begin{align*}
U(T,V)
=
cRNT-\frac{aN^2}{V}
\end{align*}

と書ける。

よって、

\begin{align*}
\left(
\frac{\partial U}{\partial V}
\right)_T
=
\frac{aN^2}{V^2}
\end{align*}

を得る。

\subsection{(iii)}

断熱過程ではエントロピーは一定である。

したがって、

\begin{align*}
S=S_0
\end{align*}

とみなせるので、

\begin{align*}
U+\frac{aN^2}{V}
=
U_0
\left(
\frac{V_0}{V-bN}
\right)^{1/c}
\exp\left(
\frac{S-S_0}{cRN}
\right)
\end{align*}

から、

\begin{align*}
U+\frac{aN^2}{V}
=
K(V-bN)^{-1/c}
\end{align*}

と書ける。ただし

\begin{align*}
K
=
U_0V_0^{1/c}
\exp\left(
\frac{S-S_0}{cRN}
\right)
\end{align*}

は断熱過程中一定である。

したがって初期状態と最終状態について

\begin{align*}
U_1+\frac{aN^2}{V_1}
=
K(V_1-bN)^{-1/c}
\end{align*}

および

\begin{align*}
U_2+\frac{aN^2}{V_2}
=
K(V_2-bN)^{-1/c}
\end{align*}

が成り立つ。

両式の比を取ると、

\begin{align*}
\frac{
U_2+\dfrac{aN^2}{V_2}
}{
U_1+\dfrac{aN^2}{V_1}
}
=
\left(
\frac{V_1-bN}{V_2-bN}
\right)^{1/c}
\end{align*}

を得る。

よって、

\begin{align*}
\frac{
U_2+aN^2/V_2
}{
U_1+aN^2/V_1
}
=
\left(
\frac{V_1-bN}{V_2-bN}
\right)^{1/c}
\end{align*}

が示された。

\subsection{(iv)}

小問(i)より、

\begin{align*}
U+\frac{aN^2}{V}
=
cRNT
\end{align*}

である。

これを小問(iii)の結果へ代入すると、

\begin{align*}
\frac{T_2}{T_1}
=
\left(
\frac{V_1-bN}{V_2-bN}
\right)^{1/c}
\end{align*}

を得る。

したがって、

\begin{align*}
T(V-bN)^{1/c}
=
\mathrm{const.}
\end{align*}

である。

さらに van der Waals 状態方程式

\begin{align*}
\left(
p+\frac{aN^2}{V^2}
\right)
(V-bN)
=
NRT
\end{align*}

を用いて温度を消去すると、

\begin{align*}
\left(
p+\frac{aN^2}{V^2}
\right)
(V-bN)
(V-bN)^{1/c}
=
\mathrm{const.}
\end{align*}

となる。

すなわち、

\begin{align*}
\left(
p+\frac{aN^2}{V^2}
\right)
(V-bN)^{1+1/c}
=
\mathrm{const.}
\end{align*}

である。

ここで

\begin{align*}
\gamma
=
\frac{C_P}{C_V}
=
\frac{c+1}{c}
=
1+\frac1c
\end{align*}

を用いれば、

\begin{align*}
\left(
p+\frac{aN^2}{V^2}
\right)
(V-bN)^\gamma
=
\mathrm{const.}
\end{align*}

を得る。

以上より、van der Waals 気体に対する Poisson の関係式は

\begin{align*}
T(V-bN)^{1/c}
=
\mathrm{const.},
\qquad
\left(
p+\frac{aN^2}{V^2}
\right)
(V-bN)^\gamma
=
\mathrm{const.}
\end{align*}

である。
\section{2-6}

温度 $T$ の気体について、速度分布関数

\begin{align*}
F(v)
=
4\pi
\left(
\frac{m}{2\pi k_BT}
\right)^{3/2}
v^2
\exp\left(
-\frac{mv^2}{2k_BT}
\right)
\end{align*}

を考える。

以下では

\begin{align*}
a
=
\frac{m}{2k_BT}
\end{align*}

と置く。

\subsection{(i)}

まず分布関数が規格化されていることを確認する。

積分値は

\begin{align*}
\int_0^\infty F(v)\,dv
=
4\pi
\left(
\frac{m}{2\pi k_BT}
\right)^{3/2}
\int_0^\infty
v^2e^{-av^2}\,dv
\end{align*}

である。

ガウス積分

\begin{align*}
\int_0^\infty e^{-ax^2}\,dx
=
\frac12\sqrt{\frac{\pi}{a}}
\end{align*}

を $a$ で微分すると、

\begin{align*}
\int_0^\infty x^2e^{-ax^2}\,dx
=
\frac{\sqrt{\pi}}{4a^{3/2}}
\end{align*}

を得る。

したがって、

\begin{align*}
\int_0^\infty F(v)\,dv
&=
4\pi
\left(
\frac{m}{2\pi k_BT}
\right)^{3/2}
\frac{\sqrt{\pi}}{4a^{3/2}} \\
&=
4\pi
\left(
\frac{m}{2\pi k_BT}
\right)^{3/2}
\frac{\sqrt{\pi}}{4}
\left(
\frac{2k_BT}{m}
\right)^{3/2} \\
&=
1。
\end{align*}

よって、

\begin{align*}
\int_0^\infty F(v)\,dv=1
\end{align*}

であり、確かに規格化されている。

\subsection{(ii)}

最確速度 $v_m$ は

\begin{align*}
\frac{dF}{dv}=0
\end{align*}

を満たす点で与えられる。

指数関数部分は常に正であるから、

\begin{align*}
f(v)=v^2e^{-av^2}
\end{align*}

を微分すればよい。

\begin{align*}
\frac{df}{dv}
&=
2ve^{-av^2}
-2av^3e^{-av^2} \\
&=
2ve^{-av^2}(1-av^2)。
\end{align*}

したがって、

\begin{align*}
v=0
\end{align*}

または

\begin{align*}
1-av^2=0
\end{align*}

である。

$F(0)=0$ であるから極大値は

\begin{align*}
v_m^2=\frac1a
\end{align*}

で与えられる。

よって、

\begin{align*}
v_m
=
\sqrt{\frac{2k_BT}{m}}
\end{align*}

を得る。

\subsection{(iii)}

平均速度

\begin{align*}
\bar v
=
\int_0^\infty vF(v)\,dv
\end{align*}

を求める。

定義より、

\begin{align*}
\bar v
=
4\pi
\left(
\frac{m}{2\pi k_BT}
\right)^{3/2}
\int_0^\infty
v^3e^{-av^2}dv。
\end{align*}

ここで

\begin{align*}
u=av^2
\end{align*}

と置くと、

\begin{align*}
du=2av\,dv
\end{align*}

であり、

\begin{align*}
v^3dv
=
\frac{u}{2a^2}\,du
\end{align*}

となる。

したがって、

\begin{align*}
\int_0^\infty
v^3e^{-av^2}dv
&=
\frac1{2a^2}
\int_0^\infty ue^{-u}du \\
&=
\frac1{2a^2}
\Gamma(2) \\
&=
\frac1{2a^2}
\end{align*}

を得る。

よって、

\begin{align*}
\bar v
&=
4\pi
\left(
\frac{m}{2\pi k_BT}
\right)^{3/2}
\frac1{2}
\left(
\frac{2k_BT}{m}
\right)^2 \\
&=
\sqrt{\frac{8k_BT}{\pi m}}。
\end{align*}

したがって、

\begin{align*}
\bar v
=
\sqrt{\frac{8k_BT}{\pi m}}
\end{align*}

を得る。

\subsection{(iv)}

平均運動エネルギー

\begin{align*}
\bar E
=
\int_0^\infty
F(v)\frac12 mv^2\,dv
\end{align*}

を求める。

まず

\begin{align*}
I
=
\int_0^\infty
v^4e^{-av^2}dv
\end{align*}

を計算する。

前問と同様に

\begin{align*}
u=av^2
\end{align*}

を用いると、

\begin{align*}
I
=
\frac1{2a^{5/2}}
\int_0^\infty
u^{3/2}e^{-u}du
=
\frac1{2a^{5/2}}
\Gamma\left(\frac52\right)。
\end{align*}

ここで

\begin{align*}
\Gamma\left(\frac52\right)
=
\frac32\frac12\sqrt{\pi}
=
\frac{3\sqrt{\pi}}4
\end{align*}

であるから、

\begin{align*}
I
=
\frac{3\sqrt{\pi}}
{8a^{5/2}}
\end{align*}

となる。

したがって、

\begin{align*}
\bar E
&=
\frac12m
4\pi
\left(
\frac{m}{2\pi k_BT}
\right)^{3/2}
\frac{3\sqrt{\pi}}
{8a^{5/2}} \\
&=
\frac32k_BT。
\end{align*}

よって、

\begin{align*}
\bar E
=
\frac32k_BT
\end{align*}

を得る。

以上より、単原子理想気体の平均運動エネルギーは温度のみに比例し、

\begin{align*}
\bar E
=
\frac32k_BT
\end{align*}

であることが示された。これはエネルギー等分配則と一致する。
\section{3-1}

Carnotサイクルについて考える。作業物質は比熱比

\begin{align*}
\gamma=\frac53
\end{align*}

の単原子理想気体であり、すべての過程は準静的に行われるものとする。 【1-c217e7】

\subsection{(i)}

状態1から状態2への変化は温度

\begin{align*}
T=T_H
\end{align*}

での等温膨張である。

理想気体の状態方程式

\begin{align*}
pV=NRT_H
\end{align*}

が成り立つので、

\begin{align*}
p_1V_1
=
p_2V_2
\end{align*}

である。

問題文より

\begin{align*}
V_2=2V_1
\end{align*}

であるから、

\begin{align*}
\frac{p_2}{p_1}
=
\frac{V_1}{V_2}
=
\frac12
\end{align*}

を得る。

したがって、

\begin{align*}
\frac{p_2}{p_1}
=
\frac12
\end{align*}

である。 【1-c217e7】

\subsection{(ii)}

状態2から状態3への変化は断熱膨張である。

Poissonの関係式

\begin{align*}
pV^\gamma
=
\mathrm{const.}
\end{align*}

より、

\begin{align*}
p_2V_2^\gamma
=
p_3V_3^\gamma
\end{align*}

である。

したがって、

\begin{align*}
\frac{p_3}{p_2}
=
\left(
\frac{V_2}{V_3}
\right)^\gamma
\end{align*}

を得る。

一方、

\begin{align*}
TV^{\gamma-1}
=
\mathrm{const.}
\end{align*}

より、

\begin{align*}
T_HV_2^{\gamma-1}
=
T_LV_3^{\gamma-1}
\end{align*}

である。

したがって、

\begin{align*}
\frac{V_3}{V_2}
=
\left(
\frac{T_H}{T_L}
\right)^{1/(\gamma-1)}
\end{align*}

となる。

ここで

\begin{align*}
\gamma-1=\frac23
\end{align*}

であるので、

\begin{align*}
\frac{V_3}{V_2}
=
\left(
\frac{T_H}{T_L}
\right)^{3/2}
\end{align*}

を得る。

また、

\begin{align*}
\frac{p_3}{p_2}
&=
\left(
\frac{V_2}{V_3}
\right)^\gamma \\
&=
\left[
\left(
\frac{T_L}{T_H}
\right)^{3/2}
\right]^{5/3} \\
&=
\left(
\frac{T_L}{T_H}
\right)^{5/2}。
\end{align*}

よって、

\begin{align*}
\frac{V_3}{V_2}
=
\left(
\frac{T_H}{T_L}
\right)^{3/2},
\qquad
\frac{p_3}{p_2}
=
\left(
\frac{T_L}{T_H}
\right)^{5/2}
\end{align*}

である。 【1-c217e7】

\subsection{(iii)}

状態3から状態4への変化は温度

\begin{align*}
T=T_L
\end{align*}

での等温圧縮である。

したがって、

\begin{align*}
p_3V_3
=
p_4V_4
\end{align*}

が成り立つ。

よって、

\begin{align*}
\frac{p_4}{p_3}
=
\frac{V_3}{V_4}
\end{align*}

である。

問題文より

\begin{align*}
V_4=V_2
\end{align*}

であるから、

\begin{align*}
\frac{p_4}{p_3}
=
\frac{V_3}{V_2}
=
\left(
\frac{T_H}{T_L}
\right)^{3/2}
\end{align*}

を得る。

したがって、

\begin{align*}
\frac{p_4}{p_3}
=
\left(
\frac{T_H}{T_L}
\right)^{3/2}
\end{align*}

である。 【1-c217e7】

\subsection{(iv)}

状態4から状態1への変化は断熱圧縮である。

Poissonの関係式より、

\begin{align*}
T_LV_4^{\gamma-1}
=
T_HV_1^{\gamma-1}
\end{align*}

である。

問題文より

\begin{align*}
V_4=2V_1
\end{align*}

であるので、

\begin{align*}
\frac{T_H}{T_L}
=
\left(
\frac{V_4}{V_1}
\right)^{\gamma-1}
=
2^{\gamma-1}
\end{align*}

となる。

\begin{align*}
\gamma-1
=
\frac23
\end{align*}

を代入すると、

\begin{align*}
\frac{T_H}{T_L}
=
2^{2/3}
\end{align*}

を得る。

すなわち、

\begin{align*}
\frac{T_H}{T_L}
=
2^{2/3}
\simeq
1.587
\end{align*}

である。 【1-c217e7】

\subsection{(v)}

各状態の位置関係は次のようになる。

\begin{itemize}
\item 状態1から状態2：高温等温膨張
\item 状態2から状態3：断熱膨張
\item 状態3から状態4：低温等温圧縮
\item 状態4から状態1：断熱圧縮
\end{itemize}

また、

\begin{align*}
V_1
<
V_2=V_4
<
V_3
\end{align*}

であり、

\begin{align*}
T_H>T_L
\end{align*}

である。

したがって、$p-V$ 図では高温側の等温線が上側、低温側の等温線が下側に位置し、それらを二本の断熱曲線が接続する閉曲線となる。

概形は次の通りである。

\begin{center}
\begin{tikzpicture}[scale=1.2]
\draw[->] (0,0) -- (5,0) node[right] {$V$};
\draw[->] (0,0) -- (0,4) node[above] {$p$};

\node[left] at (1.0,3.0) {1};
\node[above] at (2.5,2.4) {2};
\node[right] at (4.0,1.2) {3};
\node[above] at (2.5,1.0) {4};

\draw[smooth] (1.0,3.0) .. controls (1.6,2.8) and (2.0,2.5) .. (2.5,2.4);
\draw[smooth] (2.5,2.4) .. controls (3.2,2.0) and (3.7,1.6) .. (4.0,1.2);
\draw[smooth] (4.0,1.2) .. controls (3.5,1.1) and (3.0,1.0) .. (2.5,1.0);
\draw[smooth] (2.5,1.0) .. controls (1.8,1.4) and (1.2,2.1) .. (1.0,3.0);

\node at (3.8,2.7) {$T=T_H$};
\node at (3.7,0.7) {$T=T_L$};
\end{tikzpicture}
\captionof*{figure}{Carnotサイクルの $p-V$ 図の概形}
\end{center}

以上より、Carnotサイクルの各状態量の比および軌跡が求められた。
\section{3-2}

オットーサイクルの効率を求める。作業物質は理想気体であり、

\begin{align*}
A\rightarrow B
\end{align*}

および

\begin{align*}
C\rightarrow D
\end{align*}

は断熱過程、

\begin{align*}
B\rightarrow C
\end{align*}

および

\begin{align*}
D\rightarrow A
\end{align*}

は定積過程である。 【1-539e2a】

\subsection{(i)}

オットーサイクルでは、

\begin{align*}
A\rightarrow B
\end{align*}

で断熱圧縮、

\begin{align*}
B\rightarrow C
\end{align*}

で定積加熱、

\begin{align*}
C\rightarrow D
\end{align*}

で断熱膨張、

\begin{align*}
D\rightarrow A
\end{align*}

で定積冷却が起こる。

したがって、$p-V$ 図では二本の断熱曲線を二本の鉛直線（定積過程）が結ぶ形となる。 【1-539e2a】

\begin{center}
\begin{tikzpicture}[scale=1.1]
\draw[->] (0,0) -- (5,0) node[right] {$V$};
\draw[->] (0,0) -- (0,4.2) node[above] {$p$};

\node[left] at (4.0,0.9) {A};
\node[left] at (2.0,2.2) {B};
\node[left] at (2.0,3.5) {C};
\node[right] at (4.0,1.5) {D};

\draw[smooth] (4.0,0.9) .. controls (3.3,1.1) and (2.5,1.7) .. (2.0,2.2);
\draw (2.0,2.2)--(2.0,3.5);
\draw[smooth] (2.0,3.5) .. controls (2.7,3.0) and (3.5,2.0) .. (4.0,1.5);
\draw (4.0,1.5)--(4.0,0.9);
\end{tikzpicture}
\captionof*{figure}{オットーサイクルの概形}
\end{center}

\subsection{(ii)}

過程

\begin{align*}
B\rightarrow C
\end{align*}

は定積過程である。

したがって、

\begin{align*}
W=0
\end{align*}

であり、熱力学第1法則より

\begin{align*}
Q
=
U_C-U_B
\end{align*}

を得る。

理想気体では

\begin{align*}
U=C_VT
\end{align*}

であるから、

\begin{align*}
Q
=
C_V(T_C-T_B)
\end{align*}

となる。

同様に、

\begin{align*}
D\rightarrow A
\end{align*}

も定積過程であるので、

\begin{align*}
Q'
=
C_V(T_D-T_A)
\end{align*}

を得る。

ここで問題文の定義では $Q'$ は外部へ放出される熱量であるから正の量として扱う。

したがって、

\begin{align*}
Q'
=
C_V(T_D-T_A)
\end{align*}

である。 【1-539e2a】

\subsection{(iii)}

断熱過程では

\begin{align*}
Q=0
\end{align*}

である。

したがって熱力学第1法則より、

\begin{align*}
W
=
U_B-U_A
=
C_V(T_B-T_A)
\end{align*}

を得る。

また、

\begin{align*}
C\rightarrow D
\end{align*}

では気体が外部へ仕事をするので、

\begin{align*}
W'
=
U_C-U_D
=
C_V(T_C-T_D)
\end{align*}

となる。

よって、

\begin{align*}
W=C_V(T_B-T_A),
\qquad
W'=C_V(T_C-T_D)
\end{align*}

である。 【1-539e2a】

\subsection{(iv)}

前問より、

\begin{align*}
W+Q-W'-Q'
&=
C_V(T_B-T_A)
+
C_V(T_C-T_B) \\
&\quad
-C_V(T_C-T_D)
-C_V(T_D-T_A) \\
&=
0。
\end{align*}

したがって、

\begin{align*}
W+Q-W'-Q'=0
\end{align*}

となり、サイクル全体で熱力学第1法則が成立することが確認できた。 【1-539e2a】

\subsection{(v)}

効率は

\begin{align*}
\eta
=
\frac{W'-W}{Q}
\end{align*}

で定義される。

小問(iv)の結果

\begin{align*}
W'-W
=
Q-Q'
\end{align*}

を用いると、

\begin{align*}
\eta
=
1-\frac{Q'}{Q}
\end{align*}

となる。

ここで、

\begin{align*}
Q
=
C_V(T_C-T_B),
\qquad
Q'
=
C_V(T_D-T_A)
\end{align*}

であるから、

\begin{align*}
\eta
=
1-
\frac{T_D-T_A}{T_C-T_B}
\end{align*}

を得る。

次に断熱過程に対する Poisson の関係式

\begin{align*}
TV^{\gamma-1}
=
\mathrm{const.}
\end{align*}

を用いる。

過程

\begin{align*}
A\rightarrow B
\end{align*}

に対して、

\begin{align*}
T_AV_2^{\gamma-1}
=
T_BV_1^{\gamma-1}
\end{align*}

であるから、

\begin{align*}
T_B
=
T_A
\left(
\frac{V_2}{V_1}
\right)^{\gamma-1}
\end{align*}

を得る。

圧縮比を

\begin{align*}
\epsilon
=
\frac{V_2}{V_1}
\end{align*}

と定義すると、

\begin{align*}
T_B=T_A\epsilon^{\gamma-1}
\end{align*}

となる。

同様に、

\begin{align*}
C\rightarrow D
\end{align*}

について

\begin{align*}
T_D
=
\frac{T_C}{\epsilon^{\gamma-1}}
\end{align*}

を得る。

したがって、

\begin{align*}
\frac{T_D-T_A}
{T_C-T_B}
&=
\frac{
\dfrac{T_C}{\epsilon^{\gamma-1}}
-\dfrac{T_B}{\epsilon^{\gamma-1}}
}
{T_C-T_B} \\
&=
\frac{1}{\epsilon^{\gamma-1}}。
\end{align*}

これを効率の式へ代入すると、

\begin{align*}
\eta
&=
1-
\frac{1}{\epsilon^{\gamma-1}}。
\end{align*}

以上より、

\begin{align*}
\eta
=
1-
\frac{1}{\epsilon^{\gamma-1}}
\end{align*}

が示された。
\section{3-3}

粒子数を含めた理想気体の内部エネルギー

\begin{align*}
U(S,V,N)
=
U_0
\left(
\frac{N}{N_0}
\right)^{\gamma}
\left(
\frac{V_0}{V}
\right)^{\gamma-1}
\exp\left[
\frac{\gamma-1}{R}
\left(
\frac{S}{N}-\frac{S_0}{N_0}
\right)
\right]
\end{align*}

を考える。ここで $\gamma=C_P/C_V$ は比熱比であり、

\begin{align*}
\gamma>1
\end{align*}

とする。 【1-e2a27e】

\subsection{(i)}

まず、内部エネルギーが示量的であることを示す。

$S,V,N$ を同時に $\lambda$ 倍すると、

\begin{align*}
U(\lambda S,\lambda V,\lambda N)
=
U_0
\left(
\frac{\lambda N}{N_0}
\right)^\gamma
\left(
\frac{V_0}{\lambda V}
\right)^{\gamma-1}
\exp\left[
\frac{\gamma-1}{R}
\left(
\frac{\lambda S}{\lambda N}
-\frac{S_0}{N_0}
\right)
\right]
\end{align*}

となる。

指数関数中では

\begin{align*}
\frac{\lambda S}{\lambda N}
=
\frac{S}{N}
\end{align*}

であるから、

\begin{align*}
U(\lambda S,\lambda V,\lambda N)
&=
U_0
\lambda^\gamma
\lambda^{-(\gamma-1)}
\left(
\frac{N}{N_0}
\right)^\gamma
\left(
\frac{V_0}{V}
\right)^{\gamma-1}
\exp\left[
\frac{\gamma-1}{R}
\left(
\frac{S}{N}
-\frac{S_0}{N_0}
\right)
\right] \\
&=
\lambda
U(S,V,N)。
\end{align*}

したがって、

\begin{align*}
U(\lambda S,\lambda V,\lambda N)
=
\lambda U(S,V,N)
\end{align*}

が成り立つ。

以上より、内部エネルギーは示量的である。 【1-e2a27e】

\subsection{(ii)}

次に、$N$ に関する凸性を調べる。

まず

\begin{align*}
A
=
U_0
\left(
\frac{V_0}{V}
\right)^{\gamma-1}
\exp\left(
-\frac{\gamma-1}{R}
\frac{S_0}{N_0}
\right)
\end{align*}

を定数とみなす。

すると、

\begin{align*}
U
=
A
N^\gamma
\exp\left(
\frac{\gamma-1}{R}\frac{S}{N}
\right)
\end{align*}

と書ける。

さらに

\begin{align*}
a=\frac{\gamma-1}{R}S
\end{align*}

と置くと、

\begin{align*}
U
=
AN^\gamma e^{a/N}
\end{align*}

となる。

まず一階微分を求める：

\begin{align*}
\frac{\partial U}{\partial N}
&=
A e^{a/N}
\left(
\gamma N^{\gamma-1}
-aN^{\gamma-2}
\right) \\
&=
Ae^{a/N}
N^{\gamma-2}
(\gamma N-a)。
\end{align*}

さらにもう一度微分すると、

\begin{align*}
\frac{\partial^2U}{\partial N^2}
=
Ae^{a/N}
N^{\gamma-4}
\left[
\gamma(\gamma-1)N^2
-2a(\gamma-1)N
+a^2
\right]。
\end{align*}

括弧内を整理すると、

\begin{align*}
\gamma(\gamma-1)N^2
-2a(\gamma-1)N
+a^2
&=
(\gamma-1)^2N^2
+
(\gamma-1)N^2
-2a(\gamma-1)N
+a^2 \\
&=
(\gamma-1)(N-a)^2
+a^2-(\gamma-1)a^2。
\end{align*}

より直接判定してもよいが、より簡単には

\begin{align*}
\gamma(\gamma-1)N^2
-2a(\gamma-1)N
+a^2
\end{align*}

は $N$ に関する二次式であり、その判別式は

\begin{align*}
D
&=
4a^2(\gamma-1)^2
-
4\gamma(\gamma-1)a^2 \\
&=
-4a^2(\gamma-1)
<
0
\end{align*}

である。

さらに、

\begin{align*}
\gamma(\gamma-1)>0
\end{align*}

であるから、この二次式は常に正である。

したがって、

\begin{align*}
\frac{\partial^2U}{\partial N^2}
>
0
\end{align*}

を得る。

よって、

\begin{align*}
U(S,V,N)
\end{align*}

は $N$ に関して下に凸である。 【1-e2a27e】

\subsection{(iii)}

まず $S$ に関する凸性を調べる。

内部エネルギーを

\begin{align*}
U
=
B
\exp\left(
\frac{\gamma-1}{RN}S
\right)
\end{align*}

と書く。ただし

\begin{align*}
B
=
U_0
\left(
\frac{N}{N_0}
\right)^\gamma
\left(
\frac{V_0}{V}
\right)^{\gamma-1}
\exp\left(
-\frac{\gamma-1}{R}
\frac{S_0}{N_0}
\right)
\end{align*}

である。

したがって、

\begin{align*}
\frac{\partial U}{\partial S}
=
\frac{\gamma-1}{RN}U
\end{align*}

であり、

\begin{align*}
\frac{\partial^2U}{\partial S^2}
=
\left(
\frac{\gamma-1}{RN}
\right)^2
U
>
0
\end{align*}

を得る。

よって、

\begin{align*}
U(S,V,N)
\end{align*}

は $S$ に関して下に凸である。

次に $V$ に関する凸性を調べる。

\begin{align*}
U
=
C
V^{-(\gamma-1)}
\end{align*}

と書けば、

\begin{align*}
\frac{\partial U}{\partial V}
=
-(\gamma-1)
CV^{-\gamma}
\end{align*}

であり、

\begin{align*}
\frac{\partial^2U}{\partial V^2}
=
\gamma(\gamma-1)
CV^{-(\gamma+1)}
\end{align*}

となる。

ここで

\begin{align*}
\gamma>1
\end{align*}

かつ

\begin{align*}
C>0
\end{align*}

であるから、

\begin{align*}
\frac{\partial^2U}{\partial V^2}
>
0
\end{align*}

を得る。

したがって、

\begin{align*}
U(S,V,N)
\end{align*}

は $V$ に関しても下に凸である。

以上より、

\begin{align*}
\frac{\partial^2U}{\partial S^2}>0,
\qquad
\frac{\partial^2U}{\partial V^2}>0,
\qquad
\frac{\partial^2U}{\partial N^2}>0
\end{align*}

であり、内部エネルギーは $S$、$V$、$N$ のいずれに対しても下に凸であることが示された。
\section{3-4}

理想気体の内部エネルギーの基本式

\begin{align*}
U(S,V,N)
=
U_0
\left(
\frac{N}{N_0}
\right)^\gamma
\left(
\frac{V_0}{V}
\right)^{\gamma-1}
\exp\left[
\frac{\gamma-1}{R}
\left(
\frac{S}{N}
-\frac{S_0}{N_0}
\right)
\right]
\end{align*}

を逆に解き、エントロピーの基本式を求める。 【1-233108】

\subsection{(i)}

まず両辺の対数を取る：

\begin{align*}
\log\frac{U}{U_0}
&=
\gamma\log\frac{N}{N_0}
+
(\gamma-1)\log\frac{V_0}{V}
+
\frac{\gamma-1}{R}
\left(
\frac{S}{N}
-
\frac{S_0}{N_0}
\right)。
\end{align*}

これを $S$ について解く。

\begin{align*}
\frac{\gamma-1}{R}
\left(
\frac{S}{N}
-
\frac{S_0}{N_0}
\right)
&=
\log\frac{U}{U_0}
-\gamma\log\frac{N}{N_0}
-(\gamma-1)\log\frac{V_0}{V}。
\end{align*}

したがって、

\begin{align*}
\frac{S}{N}
=
\frac{S_0}{N_0}
+
\frac{R}{\gamma-1}
\log\frac{U}{U_0}
-
\frac{\gamma R}{\gamma-1}
\log\frac{N}{N_0}
+
R\log\frac{V}{V_0}。
\end{align*}

両辺に $N$ を掛けると、

\begin{align*}
S(U,V,N)
=
N\frac{S_0}{N_0}
+
\frac{RN}{\gamma-1}
\log\frac{U}{U_0}
-
\frac{\gamma RN}{\gamma-1}
\log\frac{N}{N_0}
+
RN\log\frac{V}{V_0}
\end{align*}

を得る。

以上より、

\begin{align*}
S(U,V,N)
=
N\frac{S_0}{N_0}
+
\frac{RN}{\gamma-1}
\log\frac{U}{U_0}
-
\frac{\gamma RN}{\gamma-1}
\log\frac{N}{N_0}
+
RN\log\frac{V}{V_0}
\end{align*}

が理想気体のエントロピーの基本式である。 【1-233108】

\subsection{(ii)}

まず示量性について調べる。

得られた基本式において

\begin{align*}
U\rightarrow \lambda U,
\qquad
V\rightarrow \lambda V,
\qquad
N\rightarrow \lambda N
\end{align*}

とすると、

\begin{align*}
S(\lambda U,\lambda V,\lambda N)
&=
\lambda N\frac{S_0}{N_0}
+
\frac{R\lambda N}{\gamma-1}
\log\frac{\lambda U}{U_0}
\\
&\quad
-
\frac{\gamma R\lambda N}{\gamma-1}
\log\frac{\lambda N}{N_0}
+
R\lambda N
\log\frac{\lambda V}{V_0}。
\end{align*}

右辺を整理すると、

\begin{align*}
S(\lambda U,\lambda V,\lambda N)
&=
\lambda S(U,V,N)
\\
&\quad
+\lambda RN
\left[
\frac{1}{\gamma-1}
-\frac{\gamma}{\gamma-1}
+1
\right]
\log\lambda。
\end{align*}

ところが、

\begin{align*}
\frac{1}{\gamma-1}
-\frac{\gamma}{\gamma-1}
+1
=
0
\end{align*}

である。

したがって、

\begin{align*}
S(\lambda U,\lambda V,\lambda N)
=
\lambda S(U,V,N)
\end{align*}

を得る。

以上より、エントロピーは示量的である。 【1-233108】

次に凸性を調べる。

まず $U$ に関して二階微分する。

\begin{align*}
\frac{\partial S}{\partial U}
=
\frac{RN}{(\gamma-1)U}
\end{align*}

であるから、

\begin{align*}
\frac{\partial^2S}{\partial U^2}
=
-\frac{RN}{(\gamma-1)U^2}
<
0
\end{align*}

を得る。

したがって、$S$ は $U$ に関して上に凸である。

次に $V$ に関しては、

\begin{align*}
\frac{\partial S}{\partial V}
=
\frac{RN}{V}
\end{align*}

より、

\begin{align*}
\frac{\partial^2S}{\partial V^2}
=
-\frac{RN}{V^2}
<
0
\end{align*}

となる。

よって、$S$ は $V$ に関しても上に凸である。

最後に $N$ に関する凸性を調べる。

エントロピーを

\begin{align*}
S
=
AN
+
BN\log U
+
CN\log V
-
DN\log N
\end{align*}

の形に整理すると、

\begin{align*}
D
=
\frac{\gamma R}{\gamma-1}
\end{align*}

である。

したがって、

\begin{align*}
\frac{\partial S}{\partial N}
=
A+B\log U+C\log V-D(\log N+1)
\end{align*}

であり、

\begin{align*}
\frac{\partial^2S}{\partial N^2}
=
-\frac{D}{N}
=
-\frac{\gamma R}{(\gamma-1)N}
<
0
\end{align*}

を得る。

よって、$S$ は $N$ に関しても上に凸である。

以上より、

\begin{align*}
\frac{\partial^2S}{\partial U^2}<0,
\qquad
\frac{\partial^2S}{\partial V^2}<0,
\qquad
\frac{\partial^2S}{\partial N^2}<0
\end{align*}

であり、エントロピーは示量的であり、さらに $U$、$V$、$N$ の各変数に対して上に凸な関数であることが示された。
\section{3-5}

下に凸で微分可能な関数

\begin{align*}
f(x)
\end{align*}

に対して、

\begin{align*}
g(p)
=
\inf_x\{f(x)-px\}
\end{align*}

を Legendre 変換という。ここでは具体例と一般論の両方を調べる。 【1-4ec1a7】

\subsection{(i-a)}

\begin{align*}
f(x)=x^2+2x-1
\end{align*}

とする。

Legendre変換は

\begin{align*}
g(p)
=
\inf_x
\left(
x^2+2x-1-px
\right)
\end{align*}

である。

括弧内を

\begin{align*}
\phi(x)
=
x^2+(2-p)x-1
\end{align*}

とおく。

極値条件

\begin{align*}
\frac{d\phi}{dx}
=
2x+(2-p)
=
0
\end{align*}

より、

\begin{align*}
x_0
=
\frac{p-2}{2}
\end{align*}

を得る。

\begin{align*}
\frac{d^2\phi}{dx^2}=2>0
\end{align*}

であるから、これは最小値を与える。

したがって、

\begin{align*}
g(p)
&=
f(x_0)-px_0 \\
&=
x_0^2+(2-p)x_0-1 \\
&=
\left(\frac{p-2}{2}\right)^2
-(p-2)\frac{p-2}{2}
-1 \\
&=
-\frac{(p-2)^2}{4}-1。
\end{align*}

整理すると、

\begin{align*}
g(p)
=
-p+\frac{p^2}{4}-2
\end{align*}

ではなく、

\begin{align*}
g(p)
=
-\frac{(p-2)^2}{4}-1
\end{align*}

である。

以上より、

\begin{align*}
g(p)
=
-\frac{(p-2)^2}{4}-1
\end{align*}

を得る。 【1-4ec1a7】

\subsection{(i-b)}

求めたLegendre変換を微分すると、

\begin{align*}
g'(p)
=
-\frac12(p-2)
\end{align*}

であり、

\begin{align*}
g''(p)
=
-\frac12
<
0
\end{align*}

となる。

したがって、

\begin{align*}
g(p)
\end{align*}

は上に凸な関数である。

以上より、Legendre変換はこの例では下に凸関数を上に凸関数へ写している。 【1-4ec1a7】

\subsection{(i-c)}

接線

\begin{align*}
y=g(p)+px
\end{align*}

を求める。

まず、

\begin{align*}
p=-2
\end{align*}

のとき

\begin{align*}
g(-2)
=
-\frac{(-4)^2}{4}-1
=
-5
\end{align*}

であるから、

\begin{align*}
y=-5-2x
\end{align*}

を得る。

次に、

\begin{align*}
p=0
\end{align*}

のとき

\begin{align*}
g(0)
=-2
\end{align*}

であるから、

\begin{align*}
y=-2
\end{align*}

となる。

最後に、

\begin{align*}
p=2
\end{align*}

のとき

\begin{align*}
g(2)
=-1
\end{align*}

であるから、

\begin{align*}
y=2x-1
\end{align*}

となる。

よって図示すべき直線は

\begin{align*}
y=-5-2x,
\qquad
y=-2,
\qquad
y=2x-1
\end{align*}

である。 【1-4ec1a7】

\subsection{(i-d)}

小問(i-a)で得た

\begin{align*}
g(p)
=
-\frac{(p-2)^2}{4}-1
\end{align*}

を用いる。

\begin{align*}
g(p)+px
=
-\frac{(p-2)^2}{4}-1+px
\end{align*}

である。

これを $p$ について最大化する。

微分すると、

\begin{align*}
\frac{d}{dp}
\left[g(p)+px\right]
=
-\frac12(p-2)+x
\end{align*}

であるので、

\begin{align*}
-\frac12(p-2)+x=0
\end{align*}

より、

\begin{align*}
p=2x+2
\end{align*}

を得る。

これを代入すると、

\begin{align*}
g(p)+px
&=
-\frac{(2x)^2}{4}-1
+(2x+2)x \\
&=
-x^2-1+2x^2+2x \\
&=
x^2+2x-1。
\end{align*}

これは元の関数

\begin{align*}
f(x)
=
x^2+2x-1
\end{align*}

に一致する。

よって、

\begin{align*}
f(x)
=
\sup_p
\{g(p)+px\}
\end{align*}

が成立することが確認された。 【1-4ec1a7】

\subsection{(ii-a)}

一般の下に狭義凸な関数について考える。

\begin{align*}
g(p)
=
\inf_x
\{f(x)-px\}
\end{align*}

において、

\begin{align*}
f(x)-px
\end{align*}

が最小となる点を

\begin{align*}
x=x_0
\end{align*}

とする。

極値条件より、

\begin{align*}
\frac{d}{dx}
\left(
f(x)-px
\right)
=
f'(x)-p
=
0
\end{align*}

である。

したがって、

\begin{align*}
p=f'(x_0)
\end{align*}

を得る。

$f$ は下に狭義凸であるので、

\begin{align*}
f''(x)>0
\end{align*}

が成り立つ。

したがって、

\begin{align*}
f'(x)
\end{align*}

は単調増加であり逆関数

\begin{align*}
x_0=(f')^{-1}(p)
\end{align*}

が存在する。

さらに

\begin{align*}
p=f'(x_0)
\end{align*}

を $p$ で微分すると、

\begin{align*}
1
=
f''(x_0)
\frac{dx_0}{dp}
\end{align*}

であるから、

\begin{align*}
\frac{dx_0}{dp}
=
\frac{1}{f''(x_0)}
>
0
\end{align*}

を得る。

よって、

\begin{align*}
x_0(p)
\end{align*}

は単調増加関数である。 【1-4ec1a7】

\subsection{(ii-b)}

Legendre変換は

\begin{align*}
g(p)
=
f(x_0)-px_0
\end{align*}

と書ける。

これを微分すると、

\begin{align*}
g'(p)
&=
f'(x_0)\frac{dx_0}{dp}
-x_0
-p\frac{dx_0}{dp}。
\end{align*}

ここで

\begin{align*}
f'(x_0)=p
\end{align*}

を用いると、

\begin{align*}
g'(p)
=
-x_0
\end{align*}

となる。

さらに微分すると、

\begin{align*}
g''(p)
=
-\frac{dx_0}{dp}。
\end{align*}

前問より、

\begin{align*}
\frac{dx_0}{dp}>0
\end{align*}

であるので、

\begin{align*}
g''(p)<0
\end{align*}

を得る。

したがって、

\begin{align*}
g(p)
\end{align*}

は上に狭義凸な関数である。

以上より、下に狭義凸な関数の Legendre 変換は常に上に狭義凸となることが示された。
\section{3-6}

内部エネルギーの基本式

\begin{align*}
U(S)
=
k\exp(mS)
\end{align*}

を考える。ただし、

\begin{align*}
k
&=
U_0
\left(
\frac{N}{N_0}
\right)^\gamma
\left(
\frac{V_0}{V}
\right)^{\gamma-1}
\exp\left(
-\frac{\gamma-1}{RN_0}S_0
\right),
\\
m
&=
\frac{\gamma-1}{RN}
\end{align*}

であり、この問題では $V,N$ は定数として扱う。 【1-0a513e】

\subsection{(i)}

まず温度を求める。

熱力学の定義より、

\begin{align*}
T
=
\left(
\frac{\partial U}{\partial S}
\right)_V
\end{align*}

である。

したがって、

\begin{align*}
T
&=
km\exp(mS)
\\
&=
mU
\end{align*}

を得る。

ゆえに、

\begin{align*}
U
=
\frac{T}{m}
\end{align*}

である。

また、

\begin{align*}
T
=
km\exp(mS)
\end{align*}

を $S$ について解くと、

\begin{align*}
\exp(mS)
=
\frac{T}{km}
\end{align*}

であり、

\begin{align*}
mS
=
\log
\left(
\frac{T}{km}
\right)
\end{align*}

となる。

よって、

\begin{align*}
S(T)
=
\frac{1}{m}
\log
\left(
\frac{T}{km}
\right)
\end{align*}

を得る。 【1-0a513e】

\subsection{(ii)}

Helmholtz自由エネルギーは

\begin{align*}
F=U-TS
\end{align*}

で定義される。

前問の結果

\begin{align*}
U=\frac{T}{m}
\end{align*}

および

\begin{align*}
S
=
\frac1m
\log
\left(
\frac{T}{km}
\right)
\end{align*}

を代入すると、

\begin{align*}
F(T)
&=
\frac{T}{m}
-
\frac{T}{m}
\log
\left(
\frac{T}{km}
\right)
\\
&=
\frac{T}{m}
\left[
1-
\log
\left(
\frac{T}{km}
\right)
\right]。
\end{align*}

よって、

\begin{align*}
F(T)
=
\frac{T}{m}
\left[
1-
\log
\left(
\frac{T}{km}
\right)
\right]
\end{align*}

を得る。 【1-0a513e】

\subsection{(iii)}

今度は Legendre変換の定義

\begin{align*}
F(T)
=
\inf_S
\{U(S)-TS\}
\end{align*}

から直接求める。

まず

\begin{align*}
\Phi(S)
=
U(S)-TS
=
ke^{mS}-TS
\end{align*}

とおく。

極値条件は

\begin{align*}
\frac{d\Phi}{dS}
=
kme^{mS}-T
=
0
\end{align*}

である。

したがって、

\begin{align*}
e^{m\bar S}
=
\frac{T}{km}
\end{align*}

となる。

よって、

\begin{align*}
\bar S
=
\frac1m
\log
\left(
\frac{T}{km}
\right)
\end{align*}

を得る。

また、

\begin{align*}
\frac{d^2\Phi}{dS^2}
=
km^2e^{mS}
>
0
\end{align*}

である。

したがって、この解は確かに最小値を与える。 【1-0a513e】

\subsection{(iv)}

前問で求めた

\begin{align*}
\bar S
=
\frac1m
\log
\left(
\frac{T}{km}
\right)
\end{align*}

を用いる。

まず、

\begin{align*}
U(\bar S)
&=
ke^{m\bar S}
\\
&=
k
\frac{T}{km}
\\
&=
\frac{T}{m}
\end{align*}

である。

したがって、

\begin{align*}
F
&=
U(\bar S)-T\bar S
\\
&=
\frac{T}{m}
-
\frac{T}{m}
\log
\left(
\frac{T}{km}
\right)
\\
&=
\frac{T}{m}
\left[
1-
\log
\left(
\frac{T}{km}
\right)
\right]。
\end{align*}

よって、

\begin{align*}
F(T)
=
\frac{T}{m}
\left[
1-
\log
\left(
\frac{T}{km}
\right)
\right]
\end{align*}

を得る。 【1-0a513e】

\subsection{(v)}

小問(ii)および小問(iv)で得られた結果は

\begin{align*}
F(T)
=
\frac{T}{m}
\left[
1-
\log
\left(
\frac{T}{km}
\right)
\right]
\end{align*}

で完全に一致する。

これは

\begin{align*}
F
=
U-TS
\end{align*}

という式が、任意の $S$ に対して成立する式ではなく、

\begin{align*}
T
=
\left(
\frac{\partial U}{\partial S}
\right)_V
\end{align*}

を満たす平衡状態の $S=\bar S$ を代入した結果だからである。

熱力学では内部エネルギー

\begin{align*}
U(S)
\end{align*}

が下に凸であり、

\begin{align*}
U(S)-TS
\end{align*}

も下に凸となる。

したがって極値条件

\begin{align*}
\frac{d}{dS}
\left(
U-TS
\right)
=
0
\end{align*}

を解いて得られる点は唯一の最小点である。

そのため、

\begin{align*}
F
=
\inf_S\{U(S)-TS\}
\end{align*}

という本来の Legendre変換の定義を毎回書かなくても、

\begin{align*}
F=U-TS
\end{align*}

という熱力学の公式だけで計算できるのである。

以上より、Helmholtz自由エネルギーは内部エネルギーの Legendre変換そのものであり、熱力学で用いる

\begin{align*}
F=U-TS
\end{align*}

という定義は、Legendre変換の最小化操作を省略して書いたものと解釈できる。
\section{3-7}

高温熱源と低温熱源の間で動作するヒートポンプを考える。

低温側から吸収する熱を $Q_L$、高温側へ排出する熱を $Q_H$、外部から加える仕事を $W$ とする。

冷却器としての効率は

\begin{align*}
\eta
=
\frac{Q_L}{W}
\end{align*}

で定義される。 【1-ba048f】

\subsection{(i)}

ヒートポンプ全体に対して熱力学第1法則を適用する。

サイクル運転では作業物質の内部エネルギー変化は

\begin{align*}
\Delta U=0
\end{align*}

である。

したがって、

\begin{align*}
Q_L+W-Q_H=0
\end{align*}

となる。

よって、

\begin{align*}
Q_H
=
Q_L+W
\end{align*}

を得る。

以上より、

\begin{align*}
Q_H
=
Q_L+W
\end{align*}

である。 【1-ba048f】

\subsection{(ii)}

高温熱源の内部エネルギーを $U_H$ とする。

ヒートポンプ動作後には高温熱源は熱 $Q_H$ を受け取るので、

\begin{align*}
U_H
\rightarrow
U_H+Q_H
\end{align*}

となる。

したがってエントロピー変化は

\begin{align*}
\Delta S_H
=
S_H(U_H+Q_H)-S_H(U_H)
\end{align*}

である。

前問 3-4 で示したように、エントロピーは内部エネルギーに関して上に凸である。

したがって接線は関数の上側に位置するので、

\begin{align*}
S_H(U_H+Q_H)
-
S_H(U_H)
\le
\left(
\frac{\partial S_H}{\partial U_H}
\right)_{V,N}
Q_H
\end{align*}

が成り立つ。

熱力学の基本関係式より、

\begin{align*}
\left(
\frac{\partial S}{\partial U}
\right)_{V,N}
=
\frac1T
\end{align*}

であるから、

\begin{align*}
\Delta S_H
\le
\frac{Q_H}{T_H}
\end{align*}

を得る。

以上より、

\begin{align*}
\Delta S_H
\le
\frac{Q_H}{T_H}
\end{align*}

が示された。 【1-ba048f】

\subsection{(iii)}

低温熱源について同様に考える。

低温熱源は熱量 $Q_L$ を失うので、

\begin{align*}
U_L
\rightarrow
U_L-Q_L
\end{align*}

である。

したがって、

\begin{align*}
\Delta S_L
=
S_L(U_L-Q_L)-S_L(U_L)
\end{align*}

となる。

エントロピーの上凸性より、

\begin{align*}
\Delta S_L
\le
-\frac{Q_L}{T_L}
\end{align*}

を得る。

一方、熱力学第2法則より、

\begin{align*}
\Delta S_H+\Delta S_L
\ge
0
\end{align*}

である。

したがって、

\begin{align*}
\frac{Q_H}{T_H}
-
\frac{Q_L}{T_L}
\ge
0
\end{align*}

を得る。

よって、

\begin{align*}
\frac{Q_H}{Q_L}
\ge
\frac{T_H}{T_L}
\end{align*}

が成り立つ。

以上より、

\begin{align*}
\frac{Q_H}{Q_L}
\ge
\frac{T_H}{T_L}
\end{align*}

が示された。これは Clausius の不等式である。 【1-ba048f】

\subsection{(iv)}

小問(i)より、

\begin{align*}
W
=
Q_H-Q_L
\end{align*}

である。

したがって効率は

\begin{align*}
\eta
=
\frac{Q_L}{Q_H-Q_L}
\end{align*}

と書ける。

ここで Clausius の不等式

\begin{align*}
Q_H
\ge
Q_L\frac{T_H}{T_L}
\end{align*}

を代入すると、

\begin{align*}
W
=
Q_H-Q_L
\ge
Q_L
\left(
\frac{T_H}{T_L}-1
\right)
\end{align*}

を得る。

よって、

\begin{align*}
\eta
=
\frac{Q_L}{W}
\le
\frac{
Q_L
}{
Q_L
\left(
\dfrac{T_H}{T_L}-1
\right)
}
\end{align*}

となる。

これを整理すると、

\begin{align*}
\eta
&\le
\frac{1}{
\dfrac{T_H-T_L}{T_L}
}
\\
&=
\frac{T_L}{T_H-T_L}。
\end{align*}

したがって、

\begin{align*}
\eta
\le
\frac{T_L}{T_H-T_L}
\end{align*}

を得る。

等号が成り立つのは可逆過程の場合であり、このとき Carnot ヒートポンプに対応する。

以上より、ヒートポンプの効率は

\begin{align*}
\eta_{\max}
=
\frac{T_L}{T_H-T_L}
\end{align*}

で与えられることが示された。
\section{4-1}

内部エネルギーの示量性から Euler の関係式および Gibbs-Duhem 関係式を導く。 【1-7d7ed4】

\subsection{(i)}

内部エネルギーは示量的であるから、

\begin{align*}
U(\lambda S,\lambda V,\lambda N)
=
\lambda U(S,V,N)
\end{align*}

が成り立つ。 【1-7d7ed4】

両辺を $\lambda$ で微分すると、

\begin{align*}
\frac{d}{d\lambda}
U(\lambda S,\lambda V,\lambda N)
=
U(S,V,N)
\end{align*}

である。

左辺に連鎖律を適用すると、

\begin{align*}
\left(
\frac{\partial U}{\partial S}
\right)_{\lambda V,\lambda N}
S
+
\left(
\frac{\partial U}{\partial V}
\right)_{\lambda S,\lambda N}
V
+
\left(
\frac{\partial U}{\partial N}
\right)_{\lambda S,\lambda V}
N
=
U(S,V,N)。
\end{align*}

ここで $\lambda=1$ とおく。

さらに熱力学の定義

\begin{align*}
T
=
\left(
\frac{\partial U}{\partial S}
\right)_{V,N},
\qquad
p
=
-
\left(
\frac{\partial U}{\partial V}
\right)_{S,N},
\qquad
\mu
=
\left(
\frac{\partial U}{\partial N}
\right)_{S,V}
\end{align*}

を用いると、

\begin{align*}
TS-pV+\mu N
=
U
\end{align*}

を得る。

したがって、

\begin{align*}
U
=
TS-pV+\mu N
\end{align*}

が示された。

これを Euler の関係式という。 【1-7d7ed4】

\subsection{(ii)}

Helmholtz 自由エネルギー

\begin{align*}
F=U-TS
\end{align*}

について Euler の関係式を代入すると、

\begin{align*}
F
&=
(TS-pV+\mu N)-TS \\
&=
-pV+\mu N。
\end{align*}

したがって、

\begin{align*}
F=-pV+\mu N
\end{align*}

を得る。

次にエンタルピー

\begin{align*}
H=U+pV
\end{align*}

に対して、

\begin{align*}
H
&=
(TS-pV+\mu N)+pV \\
&=
TS+\mu N。
\end{align*}

よって、

\begin{align*}
H=TS+\mu N
\end{align*}

である。

さらに Gibbs 自由エネルギー

\begin{align*}
G
=
U-TS+pV
\end{align*}

について、

\begin{align*}
G
&=
(TS-pV+\mu N)
-TS+pV \\
&=
\mu N。
\end{align*}

したがって、

\begin{align*}
G=\mu N
\end{align*}

を得る。

以上より、

\begin{align*}
F&=-pV+\mu N, \\
H&=TS+\mu N, \\
G&=\mu N
\end{align*}

が各熱力学関数に対する Euler の関係式である。 【1-7d7ed4】

\subsection{(iii)}

内部エネルギーの Euler の関係式

\begin{align*}
U
=
TS-pV+\mu N
\end{align*}

を Legendre 変換へ代入する。

Helmholtz 自由エネルギーは

\begin{align*}
F
=
U-TS
\end{align*}

であるから、

\begin{align*}
F
=
-pV+\mu N
\end{align*}

を得る。

同様にエンタルピーは

\begin{align*}
H
=
U+pV
\end{align*}

より、

\begin{align*}
H
=
TS+\mu N
\end{align*}

となる。

さらに Gibbs 自由エネルギーは

\begin{align*}
G
=
U-TS+pV
\end{align*}

より、

\begin{align*}
G
=
\mu N
\end{align*}

である。

これらは小問(ii)で示した結果と完全に一致する。

したがって、Euler の関係式と Legendre 変換は互いに矛盾しない。 【1-7d7ed4】

\subsection{(iv)}

Euler の関係式

\begin{align*}
U
=
TS-pV+\mu N
\end{align*}

を微分する：

\begin{align*}
dU
=
TdS+SdT-p\,dV-V\,dp+\mu\,dN+N\,d\mu。
\end{align*}

一方、熱力学の基本関係式より、

\begin{align*}
dU
=
TdS-p\,dV+\mu\,dN
\end{align*}

である。

両式を比較すると、

\begin{align*}
SdT-Vdp+Nd\mu
=
0
\end{align*}

を得る。

したがって、

\begin{align*}
Nd\mu+SdT-Vdp
=
0
\end{align*}

が示された。

これを Gibbs-Duhem 関係式という。
\section{4-2}

Maxwell の関係式を用いて熱容量の微分公式を導く。 【1-05740d】

\subsection{(i)}

示すべき式は

\begin{align*}
\left(
\frac{\partial C_V}{\partial V}
\right)_T
=
T
\left(
\frac{\partial^2 p}{\partial T^2}
\right)_V
\end{align*}

である。 【1-05740d】

まず Helmholtz 自由エネルギー

\begin{align*}
F=U-TS
\end{align*}

を考える。

その全微分は

\begin{align*}
dF
&=
dU-TdS-SdT \\
&=
(TdS-pdV)-TdS-SdT \\
&=
-SdT-pdV
\end{align*}

である。

したがって、

\begin{align*}
S
=
-\left(
\frac{\partial F}{\partial T}
\right)_V,
\qquad
p
=
-\left(
\frac{\partial F}{\partial V}
\right)_T
\end{align*}

を得る。

混合偏導関数の交換可能性から、

\begin{align*}
\frac{\partial}{\partial V}
\left(
\frac{\partial F}{\partial T}
\right)_V
=
\frac{\partial}{\partial T}
\left(
\frac{\partial F}{\partial V}
\right)_T
\end{align*}

であるので、

\begin{align*}
\left(
\frac{\partial S}{\partial V}
\right)_T
=
\left(
\frac{\partial p}{\partial T}
\right)_V
\end{align*}

を得る。

これが Maxwell の関係式である。

次に定積熱容量

\begin{align*}
C_V
=
T
\left(
\frac{\partial S}{\partial T}
\right)_V
\end{align*}

を用いる。

両辺を $V$ で偏微分すると、

\begin{align*}
\left(
\frac{\partial C_V}{\partial V}
\right)_T
&=
T
\frac{\partial}{\partial V}
\left[
\left(
\frac{\partial S}{\partial T}
\right)_V
\right] \\
&=
T
\frac{\partial}{\partial T}
\left[
\left(
\frac{\partial S}{\partial V}
\right)_T
\right]。
\end{align*}

ここで Maxwell 関係式を代入すると、

\begin{align*}
\left(
\frac{\partial C_V}{\partial V}
\right)_T
&=
T
\frac{\partial}{\partial T}
\left[
\left(
\frac{\partial p}{\partial T}
\right)_V
\right] \\
&=
T
\left(
\frac{\partial^2p}{\partial T^2}
\right)_V。
\end{align*}

よって、

\begin{align*}
\left(
\frac{\partial C_V}{\partial V}
\right)_T
=
T
\left(
\frac{\partial^2p}{\partial T^2}
\right)_V
\end{align*}

が示された。 【1-05740d】

\subsection{(ii)}

示すべき式は

\begin{align*}
\left(
\frac{\partial C_P}{\partial p}
\right)_T
=
-T
\left(
\frac{\partial^2V}{\partial T^2}
\right)_p
\end{align*}

である。 【1-05740d】

Gibbs 自由エネルギー

\begin{align*}
G
=
U-TS+pV
\end{align*}

を考える。

全微分は

\begin{align*}
dG
&=
dU-TdS-SdT+pdV+Vdp \\
&=
(TdS-pdV)-TdS-SdT+pdV+Vdp \\
&=
-SdT+Vdp
\end{align*}

である。

したがって、

\begin{align*}
S
=
-\left(
\frac{\partial G}{\partial T}
\right)_p,
\qquad
V
=
\left(
\frac{\partial G}{\partial p}
\right)_T
\end{align*}

を得る。

混合偏導関数の交換可能性より、

\begin{align*}
-
\left(
\frac{\partial S}{\partial p}
\right)_T
=
\left(
\frac{\partial V}{\partial T}
\right)_p
\end{align*}

すなわち、

\begin{align*}
\left(
\frac{\partial S}{\partial p}
\right)_T
=
-
\left(
\frac{\partial V}{\partial T}
\right)_p
\end{align*}

が成り立つ。

これも Maxwell の関係式である。

次に定圧熱容量

\begin{align*}
C_P
=
T
\left(
\frac{\partial S}{\partial T}
\right)_p
\end{align*}

を用いる。

両辺を圧力で偏微分すると、

\begin{align*}
\left(
\frac{\partial C_P}{\partial p}
\right)_T
&=
T
\frac{\partial}{\partial p}
\left[
\left(
\frac{\partial S}{\partial T}
\right)_p
\right] \\
&=
T
\frac{\partial}{\partial T}
\left[
\left(
\frac{\partial S}{\partial p}
\right)_T
\right]。
\end{align*}

ここで Maxwell 関係式を代入すると、

\begin{align*}
\left(
\frac{\partial C_P}{\partial p}
\right)_T
&=
-
T
\frac{\partial}{\partial T}
\left[
\left(
\frac{\partial V}{\partial T}
\right)_p
\right] \\
&=
-
T
\left(
\frac{\partial^2V}{\partial T^2}
\right)_p。
\end{align*}

したがって、

\begin{align*}
\left(
\frac{\partial C_P}{\partial p}
\right)_T
=
-T
\left(
\frac{\partial^2V}{\partial T^2}
\right)_p
\end{align*}

が示された。 
\section{4-3}

ビリアル展開

\begin{align*}
p
=
\frac{NRT}{V}
\left(
1+\frac{N}{V}B(T)
\right)
\end{align*}

および定積熱容量

\begin{align*}
C_V
=
C_V^{\mathrm{理想気体}}
-\frac{N^2R}{V}Y(T)
\end{align*}

を考える。

Maxwell の関係式を用いて、関数 $B(T)$ と $Y(T)$ が満たす関係式を導く。 【1-19b318】

\subsection{(i)}

第4-2問で導いた関係式

\begin{align*}
\left(
\frac{\partial C_V}{\partial V}
\right)_T
=
T
\left(
\frac{\partial^2p}{\partial T^2}
\right)_V
\end{align*}

を用いる。 【1-19b318】

まず左辺を計算する。

\begin{align*}
C_V
=
C_V^{\mathrm{理想気体}}
-\frac{N^2R}{V}Y(T)
\end{align*}

であり、

\begin{align*}
C_V^{\mathrm{理想気体}}
\end{align*}

は体積に依存しないので、

\begin{align*}
\left(
\frac{\partial C_V}{\partial V}
\right)_T
&=
\frac{N^2R}{V^2}Y(T)。
\end{align*}

次に右辺を計算する。

状態方程式を展開すると、

\begin{align*}
p
&=
\frac{NRT}{V}
+
\frac{N^2RT}{V^2}B(T)。
\end{align*}

したがって、

\begin{align*}
\left(
\frac{\partial p}{\partial T}
\right)_V
&=
\frac{NR}{V}
+
\frac{N^2R}{V^2}
\left[
B(T)+T B'(T)
\right]
\end{align*}

であり、

\begin{align*}
\left(
\frac{\partial^2p}{\partial T^2}
\right)_V
&=
\frac{N^2R}{V^2}
\left[
2B'(T)
+
TB''(T)
\right]。
\end{align*}

したがって、

\begin{align*}
T
\left(
\frac{\partial^2p}{\partial T^2}
\right)_V
=
\frac{N^2R}{V^2}
\left[
2T B'(T)
+
T^2 B''(T)
\right]。
\end{align*}

これを Maxwell 関係式へ代入すると、

\begin{align*}
\frac{N^2R}{V^2}Y(T)
=
\frac{N^2R}{V^2}
\left[
2TB'(T)
+
T^2B''(T)
\right]
\end{align*}

を得る。

両辺の共通因子を消去して、

\begin{align*}
Y(T)
=
2TB'(T)
+
T^2B''(T)
\end{align*}

となる。

以上より、

\begin{align*}
Y(T)
=
2T\frac{dB}{dT}
+
T^2\frac{d^2B}{dT^2}
\end{align*}

が求める関係式である。 【1-19b318】

\subsection{(ii)}

上の結果は次の形にも書ける。

\begin{align*}
Y(T)
&=
T
\frac{d}{dT}
\left[
T
\frac{dB}{dT}
\right]。
\end{align*}

実際、

\begin{align*}
T
\frac{d}{dT}
\left[
T
\frac{dB}{dT}
\right]
&=
T
\left[
\frac{dB}{dT}
+
T\frac{d^2B}{dT^2}
\right] \\
&=
2T\frac{dB}{dT}
+
T^2\frac{d^2B}{dT^2}
\end{align*}

である。

したがって、

\begin{align*}
Y(T)
=
T
\frac{d}{dT}
\left[
T
\frac{dB}{dT}
\right]
\end{align*}

とも表せる。

この式により、ビリアル係数 $B(T)$ の温度依存性が分かれば、熱容量補正 $Y(T)$ が Maxwell の関係式から決定されることが分かる。
\section{4-4}

Maxwell の関係式を用いて内部エネルギーの体積依存性を調べる。 【1-a4eeb9】

\subsection{(i)}

示すべき式は

\begin{align*}
\left(
\frac{\partial U}{\partial V}
\right)_T
=
T
\left(
\frac{\partial p}{\partial T}
\right)_V
-p
\end{align*}

である。 【1-a4eeb9】

まず熱力学第1法則

\begin{align*}
dU
=
TdS-p\,dV
\end{align*}

を用いる。

エントロピーを変数 $T,V$ の関数とみなすと、

\begin{align*}
dS
=
\left(
\frac{\partial S}{\partial T}
\right)_V dT
+
\left(
\frac{\partial S}{\partial V}
\right)_T dV
\end{align*}

である。

これを熱力学第1法則へ代入すると、

\begin{align*}
dU
&=
T
\left[
\left(
\frac{\partial S}{\partial T}
\right)_V dT
+
\left(
\frac{\partial S}{\partial V}
\right)_T dV
\right]
-p\,dV \\
&=
T
\left(
\frac{\partial S}{\partial T}
\right)_V dT
+
\left[
T
\left(
\frac{\partial S}{\partial V}
\right)_T
-p
\right]
dV。
\end{align*}

一方、

\begin{align*}
dU
=
\left(
\frac{\partial U}{\partial T}
\right)_V dT
+
\left(
\frac{\partial U}{\partial V}
\right)_T dV
\end{align*}

でもある。

したがって、$dV$ の係数を比較すると、

\begin{align*}
\left(
\frac{\partial U}{\partial V}
\right)_T
=
T
\left(
\frac{\partial S}{\partial V}
\right)_T
-p
\end{align*}

を得る。

次に Helmholtz 自由エネルギー

\begin{align*}
F=U-TS
\end{align*}

について、

\begin{align*}
dF
=
-SdT-pdV
\end{align*}

であるから、

\begin{align*}
S
=
-\left(
\frac{\partial F}{\partial T}
\right)_V,
\qquad
p
=
-\left(
\frac{\partial F}{\partial V}
\right)_T
\end{align*}

を得る。

混合偏導関数の交換可能性より、

\begin{align*}
\left(
\frac{\partial S}{\partial V}
\right)_T
=
\left(
\frac{\partial p}{\partial T}
\right)_V
\end{align*}

である。

これを代入すると、

\begin{align*}
\left(
\frac{\partial U}{\partial V}
\right)_T
=
T
\left(
\frac{\partial p}{\partial T}
\right)_V
-p
\end{align*}

が得られる。

以上より、示すべき式が証明された。 【1-a4eeb9】

\subsection{(ii)}

理想気体では状態方程式

\begin{align*}
p
=
\frac{NRT}{V}
\end{align*}

が成り立つ。 【1-a4eeb9】

したがって、

\begin{align*}
\left(
\frac{\partial p}{\partial T}
\right)_V
=
\frac{NR}{V}
\end{align*}

である。

よって、

\begin{align*}
T
\left(
\frac{\partial p}{\partial T}
\right)_V
-p
&=
T
\frac{NR}{V}
-
\frac{NRT}{V} \\
&=
0。
\end{align*}

したがって、

\begin{align*}
\left(
\frac{\partial U}{\partial V}
\right)_T
=
0
\end{align*}

を得る。

以上より、理想気体の内部エネルギーは体積に依存せず温度だけの関数であることが分かる。 【1-a4eeb9】

\subsection{(iii)}

van der Waals 気体では

\begin{align*}
\left(
p+\frac{aN^2}{V^2}
\right)
(V-bN)
=
NRT
\end{align*}

である。 【1-a4eeb9】

したがって、

\begin{align*}
p
=
\frac{NRT}{V-bN}
-\frac{aN^2}{V^2}
\end{align*}

と書ける。

よって、

\begin{align*}
\left(
\frac{\partial p}{\partial T}
\right)_V
=
\frac{NR}{V-bN}
\end{align*}

である。

これを式

\begin{align*}
\left(
\frac{\partial U}{\partial V}
\right)_T
=
T
\left(
\frac{\partial p}{\partial T}
\right)_V
-p
\end{align*}

へ代入すると、

\begin{align*}
\left(
\frac{\partial U}{\partial V}
\right)_T
&=
T
\frac{NR}{V-bN}
-
\left(
\frac{NRT}{V-bN}
-\frac{aN^2}{V^2}
\right) \\
&=
\frac{aN^2}{V^2}。
\end{align*}

したがって、

\begin{align*}
\left(
\frac{\partial U}{\partial V}
\right)_T
=
\frac{aN^2}{V^2}
\end{align*}

を得る。

これは第2回演習の van der Waals 気体の内部エネルギー

\begin{align*}
U(T,V)
=
cRNT-\frac{aN^2}{V}
\end{align*}

から直接計算した結果と一致する。

以上より、van der Waals 気体では分子間引力を表す係数 $a$ のために内部エネルギーは体積へ依存し、

\begin{align*}
\left(
\frac{\partial U}{\partial V}
\right)_T
=
\frac{aN^2}{V^2}
\end{align*}

となることが示された。
\section{4-5}

Carnot サイクルの効率に基づいて定義される熱力学的温度と、経験的温度との関係を調べる。 【1-2e53e5】

\subsection{(i)}

まず式(45)

\begin{align*}
\left(
\frac{\partial U}{\partial V}
\right)_T
=
T
\left(
\frac{\partial p}{\partial T}
\right)_V
-p
\end{align*}

を用いる。 【1-2e53e5】

熱力学的温度 $T$ の代わりに経験的温度 $\theta$ を変数として取ると、

\begin{align*}
dU
=
\left(
\frac{\partial U}{\partial \theta}
\right)_V d\theta
+
\left(
\frac{\partial U}{\partial V}
\right)_\theta dV
\end{align*}

である。

熱力学第1法則

\begin{align*}
dU=d'Q-p\,dV
\end{align*}

より、

\begin{align*}
d'Q
=
\left(
\frac{\partial U}{\partial \theta}
\right)_V d\theta
+
\left[
\left(
\frac{\partial U}{\partial V}
\right)_\theta
+p
\right]dV
\end{align*}

を得る。

したがって、等温変化に伴う熱量は

\begin{align*}
\left(
\frac{d'Q}{dV}
\right)_\theta
=
\left(
\frac{\partial U}{\partial V}
\right)_\theta
+p
\end{align*}

である。

一方、式(45)と同じ議論は任意の温度変数について成立するので、

\begin{align*}
\left(
\frac{\partial U}{\partial V}
\right)_\theta
=
T
\left(
\frac{\partial p}{\partial T}
\right)_V
-p
\end{align*}

となる。

したがって、

\begin{align*}
\left(
\frac{d'Q}{dV}
\right)_\theta
=
T
\left(
\frac{\partial p}{\partial T}
\right)_V
\end{align*}

を得る。

ここで $p=p(T(\theta),V)$ であるから、

\begin{align*}
\left(
\frac{\partial p}{\partial \theta}
\right)_V
=
\left(
\frac{\partial p}{\partial T}
\right)_V
\frac{dT}{d\theta}
\end{align*}

である。

したがって、

\begin{align*}
\left(
\frac{d'Q}{dV}
\right)_\theta
=
T
\left(
\frac{\partial p}{\partial \theta}
\right)_V
\frac{d\theta}{dT}
\end{align*}

となる。

これを整理すると、

\begin{align*}
\frac{dT}{T}
=
\frac{
\left(
\dfrac{\partial p}{\partial \theta}
\right)_V
}{
\left(
\dfrac{d'Q}{dV}
\right)_\theta
}
d\theta。
\end{align*}

基準温度 $(T_0,\theta_0)$ から積分すると、

\begin{align*}
\int_{T_0}^{T}\frac{dT'}{T'}
=
\int_{\theta_0}^{\theta}
\frac{
\left(
\dfrac{\partial p}{\partial \theta'}
\right)_V
}{
\left(
\dfrac{d'Q}{dV}
\right)_{\theta'}
}
d\theta'。
\end{align*}

よって、

\begin{align*}
\log\frac{T}{T_0}
=
\int_{\theta_0}^{\theta}
\frac{
\left(
\dfrac{\partial p}{\partial \theta'}
\right)_V
}{
\left(
\dfrac{d'Q}{dV}
\right)_{\theta'}
}
d\theta'
\end{align*}

となり、

\begin{align*}
\frac{T}{T_0}
=
\exp\left[
\int_{\theta_0}^{\theta}
\frac{
\left(
\dfrac{\partial p}{\partial \theta'}
\right)_V
}{
\left(
\dfrac{d'Q}{dV}
\right)_{\theta'}
}
d\theta'
\right]
\end{align*}

を得る。

以上より、式(46)が示された。 【1-2e53e5】

\subsection{(ii)}

経験的温度が液柱の高さに比例する場合には、

\begin{align*}
pV=k\theta
\end{align*}

が成立する。 【1-2e53e5】

したがって、

\begin{align*}
\left(
\frac{\partial p}{\partial \theta}
\right)_V
=
\frac{k}{V}
\end{align*}

である。

また Gay-Lussac-Joule の法則

\begin{align*}
\left(
\frac{\partial U}{\partial V}
\right)_\theta
=
0
\end{align*}

を用いると、

\begin{align*}
\left(
\frac{d'Q}{dV}
\right)_\theta
=
p
=
\frac{k\theta}{V}
\end{align*}

となる。

したがって、

\begin{align*}
\frac{
\left(
\dfrac{\partial p}{\partial \theta}
\right)_V
}{
\left(
\dfrac{d'Q}{dV}
\right)_\theta
}
=
\frac{k/V}{k\theta/V}
=
\frac1\theta
\end{align*}

である。

ゆえに、

\begin{align*}
\log\frac{T}{T_0}
=
\int_{\theta_0}^{\theta}
\frac{d\theta'}{\theta'}
=
\log\frac{\theta}{\theta_0}
\end{align*}

となる。

したがって、

\begin{align*}
\frac{T}{T_0}
=
\frac{\theta}{\theta_0}
\end{align*}

を得る。

よって、

\begin{align*}
T
=
\mathrm{const.}\times \theta
\end{align*}

である。

以上より、我々が実際に用いている経験的温度は熱力学的温度に比例する。 【1-2e53e5】

\subsection{(iii)}

今度は架空の歴史として、

\begin{align*}
pV=k\theta^2
\end{align*}

が経験法則として発見されていた場合を考える。 【1-2e53e5】

このとき、

\begin{align*}
\left(
\frac{\partial p}{\partial \theta}
\right)_V
=
\frac{2k\theta}{V}
\end{align*}

である。

また Gay-Lussac-Joule の法則は引き続き

\begin{align*}
\left(
\frac{\partial U}{\partial V}
\right)_\theta
=
0
\end{align*}

なので、

\begin{align*}
\left(
\frac{d'Q}{dV}
\right)_\theta
=
p
=
\frac{k\theta^2}{V}
\end{align*}

となる。

したがって、

\begin{align*}
\frac{
\left(
\dfrac{\partial p}{\partial \theta}
\right)_V
}{
\left(
\dfrac{d'Q}{dV}
\right)_\theta
}
=
\frac{2}{\theta}
\end{align*}

である。

よって、

\begin{align*}
\log\frac{T}{T_0}
&=
\int_{\theta_0}^{\theta}
\frac{2}{\theta'}\,d\theta' \\
&=
2\log\frac{\theta}{\theta_0}。
\end{align*}

したがって、

\begin{align*}
\frac{T}{T_0}
=
\left(
\frac{\theta}{\theta_0}
\right)^2
\end{align*}

を得る。

すなわち、

\begin{align*}
T
=
\mathrm{const.}\times\theta^2
\end{align*}

である。

以上より、経験的温度の定義を変えると気体法則の形は変わるが、熱力学的温度はそれに応じて再定義され、両者の間には

\begin{align*}
T\propto\theta^2
\end{align*}

という関係が成り立つ。
\section{4-6}

常磁性体に対して

\begin{align*}
dU
=
TdS+HdM
\end{align*}

が成り立つとする。ここで $H$ は外部磁場、$M$ は磁化である。 【1-9201fa】

\subsection{(i)}

まず Helmholtz 自由エネルギー

\begin{align*}
F
=
U-TS
\end{align*}

を定義する。

全微分を取ると、

\begin{align*}
dF
&=
dU-TdS-SdT \\
&=
(TdS+HdM)-TdS-SdT \\
&=
-SdT+HdM。
\end{align*}

したがって、

\begin{align*}
dF
=
-SdT+HdM
\end{align*}

を得る。

次に Gibbs 自由エネルギーを

\begin{align*}
G
=
U-TS-HM
\end{align*}

と定義する。

すると、

\begin{align*}
dG
&=
dU-TdS-SdT-HdM-MdH \\
&=
(TdS+HdM)-TdS-SdT-HdM-MdH \\
&=
-SdT-MdH。
\end{align*}

よって、

\begin{align*}
dG
=
-SdT-MdH
\end{align*}

である。

最後にエンタルピーに対応する量

\begin{align*}
H_c
=
U-HM
\end{align*}

を導入する。

その全微分は

\begin{align*}
dH_c
&=
dU-HdM-MdH \\
&=
(TdS+HdM)-HdM-MdH \\
&=
TdS-MdH。
\end{align*}

したがって、

\begin{align*}
dH_c
=
TdS-MdH
\end{align*}

を得る。

以上より、

\begin{align*}
dF&=-SdT+HdM, \\
dG&=-SdT-MdH, \\
dH_c&=TdS-MdH
\end{align*}

である。 【1-9201fa】

\subsection{(ii)}

各熱力学関数の自然変数を調べる。

内部エネルギーについては

\begin{align*}
dU
=
TdS+HdM
\end{align*}

であるから、

\begin{align*}
U=U(S,M)
\end{align*}

とみなせる。

ここで、

\begin{align*}
S,\ M
\end{align*}

はどちらも示量変数である。

また、

\begin{align*}
dH_c
=
TdS-MdH
\end{align*}

より、

\begin{align*}
H_c=H_c(S,H)
\end{align*}

であり、

\begin{align*}
S
\end{align*}

は示量変数、

\begin{align*}
H
\end{align*}

は示強変数である。

したがって、独立変数が二つとも示量変数なのは

\begin{align*}
U(S,M)
\end{align*}

のみである。

よって、

\begin{align*}
\boxed{\text{内部エネルギー }U}
\end{align*}

が該当する。 【1-9201fa】

\subsection{(iii)}

前問と同様に自然変数を見る。

\begin{align*}
dF
=
-SdT+HdM
\end{align*}

より、

\begin{align*}
F=F(T,M)
\end{align*}

である。

ここで、

\begin{align*}
T
\end{align*}

は示強変数、

\begin{align*}
M
\end{align*}

は示量変数である。

したがって両方とも示強変数ではない。

さらに、

\begin{align*}
dG
=
-SdT-MdH
\end{align*}

より、

\begin{align*}
G=G(T,H)
\end{align*}

である。

ここで

\begin{align*}
T,\ H
\end{align*}

はいずれも示強変数である。

したがって、独立変数が二つとも示強変数なのは

\begin{align*}
G(T,H)
\end{align*}

のみである。

よって、

\begin{align*}
\boxed{\text{Gibbs自由エネルギー }G}
\end{align*}

が該当する。 【1-9201fa】

\subsection{(iv)}

各熱力学関数から Maxwell 関係式を導く。

まず内部エネルギー

\begin{align*}
dU
=
TdS+HdM
\end{align*}

より、

\begin{align*}
T
=
\left(
\frac{\partial U}{\partial S}
\right)_M,
\qquad
H
=
\left(
\frac{\partial U}{\partial M}
\right)_S
\end{align*}

である。

混合偏導関数の交換可能性から、

\begin{align*}
\left(
\frac{\partial T}{\partial M}
\right)_S
=
\left(
\frac{\partial H}{\partial S}
\right)_M
\end{align*}

を得る。

次に Helmholtz 自由エネルギー

\begin{align*}
dF
=
-SdT+HdM
\end{align*}

より、

\begin{align*}
S
=
-\left(
\frac{\partial F}{\partial T}
\right)_M,
\qquad
H
=
\left(
\frac{\partial F}{\partial M}
\right)_T
\end{align*}

である。

したがって、

\begin{align*}
\left(
\frac{\partial S}{\partial M}
\right)_T
=
-
\left(
\frac{\partial H}{\partial T}
\right)_M
\end{align*}

を得る。

さらに Gibbs 自由エネルギー

\begin{align*}
dG
=
-SdT-MdH
\end{align*}

より、

\begin{align*}
S
=
-\left(
\frac{\partial G}{\partial T}
\right)_H,
\qquad
M
=
-\left(
\frac{\partial G}{\partial H}
\right)_T
\end{align*}

である。

したがって、

\begin{align*}
\left(
\frac{\partial S}{\partial H}
\right)_T
=
\left(
\frac{\partial M}{\partial T}
\right)_H
\end{align*}

を得る。

最後に

\begin{align*}
dH_c
=
TdS-MdH
\end{align*}

より、

\begin{align*}
T
=
\left(
\frac{\partial H_c}{\partial S}
\right)_H,
\qquad
M
=
-\left(
\frac{\partial H_c}{\partial H}
\right)_S
\end{align*}

である。

したがって、

\begin{align*}
\left(
\frac{\partial T}{\partial H}
\right)_S
=
-
\left(
\frac{\partial M}{\partial S}
\right)_H
\end{align*}

を得る。

以上より、4つの Maxwell 関係式

\begin{align*}
\left(
\frac{\partial T}{\partial M}
\right)_S
&=
\left(
\frac{\partial H}{\partial S}
\right)_M,
\\
\left(
\frac{\partial S}{\partial M}
\right)_T
&=
-
\left(
\frac{\partial H}{\partial T}
\right)_M,
\\
\left(
\frac{\partial S}{\partial H}
\right)_T
&=
\left(
\frac{\partial M}{\partial T}
\right)_H,
\\
\left(
\frac{\partial T}{\partial H}
\right)_S
&=
-
\left(
\frac{\partial M}{\partial S}
\right)_H
\end{align*}

が得られる。
\section{5-1}

理想常磁性体の状態方程式

\begin{align*}
M=\frac{CH}{T}
\end{align*}

を考える。ここで $C$ は Curie 定数である。 【1-214cf4】

\subsection{(i)}

示すべき式は

\begin{align*}
\left(
\frac{\partial U}{\partial M}
\right)_T
=0
\end{align*}

である。 【1-214cf4】

まず常磁性体の基本関係式

\begin{align*}
dU
=
TdS+HdM
\end{align*}

を用いる。

エントロピーを $T,M$ の関数とみなすと、

\begin{align*}
dS
=
\left(
\frac{\partial S}{\partial T}
\right)_M dT
+
\left(
\frac{\partial S}{\partial M}
\right)_T dM
\end{align*}

である。

したがって、

\begin{align*}
dU
&=
T
\left[
\left(
\frac{\partial S}{\partial T}
\right)_M dT
+
\left(
\frac{\partial S}{\partial M}
\right)_T dM
\right]
+
HdM \\
&=
T
\left(
\frac{\partial S}{\partial T}
\right)_M dT
+
\left[
T
\left(
\frac{\partial S}{\partial M}
\right)_T
+
H
\right]
dM。
\end{align*}

よって、

\begin{align*}
\left(
\frac{\partial U}{\partial M}
\right)_T
=
T
\left(
\frac{\partial S}{\partial M}
\right)_T
+
H
\end{align*}

を得る。

次に Helmholtz 自由エネルギー

\begin{align*}
F=U-TS
\end{align*}

について、

\begin{align*}
dF
=
-SdT+HdM
\end{align*}

であるから、

\begin{align*}
\left(
\frac{\partial S}{\partial M}
\right)_T
=
-
\left(
\frac{\partial H}{\partial T}
\right)_M
\end{align*}

が成り立つ。

これを代入すると、

\begin{align*}
\left(
\frac{\partial U}{\partial M}
\right)_T
=
H
-
T
\left(
\frac{\partial H}{\partial T}
\right)_M。
\end{align*}

一方、Curie の法則より、

\begin{align*}
H
=
\frac{MT}{C}
\end{align*}

である。

したがって、

\begin{align*}
\left(
\frac{\partial H}{\partial T}
\right)_M
=
\frac{M}{C}
=
\frac{H}{T}
\end{align*}

を得る。

よって、

\begin{align*}
\left(
\frac{\partial U}{\partial M}
\right)_T
&=
H-T\frac{H}{T} \\
&=
0。
\end{align*}

したがって、

\begin{align*}
\left(
\frac{\partial U}{\partial M}
\right)_T
=0
\end{align*}

が示された。

以上より、理想常磁性体の内部エネルギーは磁化 $M$ に依存せず、

\begin{align*}
U=U(T)
\end{align*}

と書ける。 【1-214cf4】

\subsection{(ii)}

今度は Curie-Weiss の法則

\begin{align*}
M
=
\frac{CH}{T-\Theta}
\end{align*}

を仮定する。ここで $\Theta$ は Curie 点である。 【1-214cf4】

この式を変形すると、

\begin{align*}
H
=
\frac{M}{C}(T-\Theta)
\end{align*}

である。

したがって、

\begin{align*}
\left(
\frac{\partial H}{\partial T}
\right)_M
=
\frac{M}{C}
\end{align*}

を得る。

前問と同様に、

\begin{align*}
\left(
\frac{\partial U}{\partial M}
\right)_T
=
H
-
T
\left(
\frac{\partial H}{\partial T}
\right)_M
\end{align*}

を用いると、

\begin{align*}
\left(
\frac{\partial U}{\partial M}
\right)_T
&=
\frac{M}{C}(T-\Theta)
-
T\frac{M}{C} \\
&=
-\frac{\Theta M}{C}。
\end{align*}

したがって、

\begin{align*}
\left(
\frac{\partial U}{\partial M}
\right)_T
=
-\frac{\Theta}{C}M
\end{align*}

を得る。

よって Curie-Weiss の法則に従う場合には内部エネルギーは磁化へ依存し、

\begin{align*}
\left(
\frac{\partial U}{\partial M}
\right)_T
\neq 0
\end{align*}

となる。

以上より、

\begin{align*}
\left(
\frac{\partial U}{\partial M}
\right)_T
=
-\frac{\Theta}{C}M
\end{align*}

が求める結果である。
\section{5-2}

理想常磁性体では Curie の法則

\begin{align*}
M=\frac{CH}{T}
\end{align*}

が成り立つものとする。また、問題5-1の結果より内部エネルギーは温度のみの関数であり、

\begin{align*}
dU=C_MdT
\end{align*}

と書ける。ただし、$C_M$ は定磁化熱容量であり、温度によらない定数とする。 【1-a786d5】

\subsection{(i)}

まず断熱曲線 $T=f(M)$ を求める。

熱力学の基本関係式は

\begin{align*}
dU
=
TdS+HdM
\end{align*}

である。

ここで

\begin{align*}
dU=C_MdT
\end{align*}

を代入すると、

\begin{align*}
TdS
=
C_MdT-HdM
\end{align*}

を得る。

断熱変化では

\begin{align*}
dS=0
\end{align*}

であるから、

\begin{align*}
C_MdT
=
HdM
\end{align*}

となる。

さらに Curie の法則

\begin{align*}
H=\frac{MT}{C}
\end{align*}

を代入すると、

\begin{align*}
C_MdT
=
\frac{MT}{C}dM
\end{align*}

である。

よって、

\begin{align*}
\frac{dT}{T}
=
\frac{M}{CC_M}dM
\end{align*}

を得る。

両辺を積分すると、

\begin{align*}
\int\frac{dT}{T}
=
\frac{1}{CC_M}
\int M\,dM
\end{align*}

より、

\begin{align*}
\log T
=
\frac{M^2}{2CC_M}
+\mathrm{const.}
\end{align*}

となる。

したがって、

\begin{align*}
T
=
A
\exp\left(
\frac{M^2}{2CC_M}
\right)
\end{align*}

を得る。ただし $A$ は積分定数である。

以上より、断熱変化では

\begin{align*}
T
=
A
\exp\left(
\frac{M^2}{2CC_M}
\right)
\end{align*}

が成り立つ。 【1-a786d5】

\subsection{(ii)}

次に磁場と磁化の関係式を求める。

Curie の法則

\begin{align*}
H=\frac{MT}{C}
\end{align*}

に前問の結果を代入すると、

\begin{align*}
H
&=
\frac{M}{C}
A
\exp\left(
\frac{M^2}{2CC_M}
\right)。
\end{align*}

したがって、

\begin{align*}
H
=
B
M
\exp\left(
\frac{M^2}{2CC_M}
\right)
\end{align*}

を得る。ただし

\begin{align*}
B=\frac{A}{C}
\end{align*}

は断熱曲線ごとに決まる定数である。

以上より、断熱変化における磁場と磁化の関係は

\begin{align*}
H
=
B
M
\exp\left(
\frac{M^2}{2CC_M}
\right)
\end{align*}

で与えられる。 【1-a786d5】

\subsection{(iii)}

前問の結果から、

\begin{align*}
T
=
A
\exp\left(
\frac{M^2}{2CC_M}
\right)
\end{align*}

である。

したがって磁化が減少すると

\begin{align*}
M^2
\end{align*}

も減少し、指数関数の値も減少する。

よって断熱的に磁化を減少させると温度は低下する。

これは断熱消磁による冷却効果に対応している。

以上より、理想常磁性体では断熱消磁によって温度が下がることが分かる。
\section{5-3}

理想常磁性体を作業物質とする Carnot サイクルを考える。

各過程は

\begin{align*}
A\to B
\end{align*}

：高温 $T_H$ の等温過程、

\begin{align*}
B\to C
\end{align*}

：断熱過程、

\begin{align*}
C\to D
\end{align*}

：低温 $T_L$ の等温過程、

\begin{align*}
D\to A
\end{align*}

：断熱過程

である。 【1-c6287d】

問題 5-2 の結果として、

\begin{align*}
T
=
A
\exp\left(
\frac{M^2}{2CC_M}
\right)
\end{align*}

が断熱曲線を与えるものとする。 【1-c6287d】

\subsection{(i)}

まず等温過程において吸収・放出される熱量を求める。

理想常磁性体では、

\begin{align*}
dU=C_MdT
\end{align*}

であるから、等温過程では

\begin{align*}
dT=0
\end{align*}

となり、

\begin{align*}
dU=0
\end{align*}

である。 【1-c6287d】

熱力学第1法則

\begin{align*}
dU=d'Q+d'W
\end{align*}

より、

\begin{align*}
d'Q=-d'W
\end{align*}

を得る。

一方、

\begin{align*}
dU
=
TdS+HdM
\end{align*}

であり、

\begin{align*}
dU=0
\end{align*}

なので、

\begin{align*}
TdS=-HdM
\end{align*}

である。

したがって、

\begin{align*}
d'Q=TdS
\end{align*}

となる。 【1-c6287d】

\subsection{(ii)}

高温等温過程 $A\to B$ において吸収される熱

\begin{align*}
Q_H
\end{align*}

を求める。

温度は一定で

\begin{align*}
T=T_H
\end{align*}

であるから、

\begin{align*}
Q_H
=
\int_A^B T_H\,dS
=
T_H(S_B-S_A)
\end{align*}

を得る。 【1-c6287d】

同様に、低温等温過程 $C\to D$ において放出される熱量は

\begin{align*}
Q_L
=
T_L(S_C-S_D)
\end{align*}

である。 【1-c6287d】

\subsection{(iii)}

断熱過程では

\begin{align*}
dS=0
\end{align*}

である。

したがって、

\begin{align*}
S_B=S_C,
\qquad
S_D=S_A
\end{align*}

が成り立つ。 【1-c6287d】

よって、

\begin{align*}
S_B-S_A
=
S_C-S_D
\end{align*}

であり、

\begin{align*}
Q_H
&=
T_H(S_B-S_A),
\\
Q_L
&=
T_L(S_B-S_A)
\end{align*}

と書ける。

したがって、

\begin{align*}
\frac{Q_L}{Q_H}
=
\frac{T_L}{T_H}
\end{align*}

を得る。 【1-c6287d】

\subsection{(iv)}

1サイクルで系が外部へする仕事を

\begin{align*}
W
=
W_{AB}
+
W_{BC}
+
W_{CD}
+
W_{DA}
\end{align*}

とする。 【1-c6287d】

サイクルでは内部エネルギーは元に戻るので、

\begin{align*}
\Delta U=0
\end{align*}

である。

したがって熱力学第1法則より、

\begin{align*}
W
=
Q_H-Q_L
\end{align*}

となる。

よって熱機関の効率

\begin{align*}
\eta
=
\frac{W}{Q_H}
\end{align*}

は、

\begin{align*}
\eta
&=
\frac{Q_H-Q_L}{Q_H}
\\
&=
1-\frac{Q_L}{Q_H}
\\
&=
1-\frac{T_L}{T_H}
\\
&=
\frac{T_H-T_L}{T_H}
\end{align*}

となる。

以上より、

\begin{align*}
\eta
=
\frac{T_H-T_L}{T_H}
\end{align*}

が示された。 【1-c6287d】

これは理想気体を作業物質とする通常の Carnot サイクルと同じ結果である。

したがって、Carnot 効率は作業物質の詳細によらず、高温熱源と低温熱源の温度だけで決まることが確認された。
\section{5-4}

磁化 $m$ を持つ磁気双極子が $N$ 個存在する系を考える。

各双極子は外部磁場 $H$ に対して平行または反平行を向くものとし、平行状態のエネルギーを

\begin{align*}
-mH
\end{align*}

反平行状態のエネルギーを

\begin{align*}
+mH
\end{align*}

とする。

また全内部エネルギーを

\begin{align*}
U=-LmH
\end{align*}

と定義する。 【1-e08515】

\subsection{(i)}

まず内部エネルギー $U$ に対応する状態数を求める。

平行な双極子の個数を $N_+$、反平行な双極子の個数を $N_-$ とすると、

\begin{align*}
N_++N_-=N
\end{align*}

である。

また、

\begin{align*}
L=N_+-N_-
\end{align*}

であるから、

\begin{align*}
N_+
&=
\frac{N+L}{2},
\\
N_-
&=
\frac{N-L}{2}
\end{align*}

を得る。

したがって、全 $N$ 個の双極子のうち

\begin{align*}
\frac{N+L}{2}
\end{align*}

個を平行状態に選ぶ組合せ数は

\begin{align*}
W(U)
&=
{}_N C_{(N+L)/2}
\\
&=
\frac{N!}
{\left(\frac{N+L}{2}\right)!
 \left(\frac{N-L}{2}\right)!}
\end{align*}

である。

よって、

\begin{align*}
W(U)
=
\frac{N!}
{\left(\frac{N+L}{2}\right)!
 \left(\frac{N-L}{2}\right)!}
\end{align*}

が示された。 【1-e08515】

\subsection{(ii)}

Boltzmann の原理より、

\begin{align*}
S(L)
=
k_B\log W(U)
\end{align*}

である。

まず Stirling の公式

\begin{align*}
\log n!
\simeq
n\log n-n
\end{align*}

を用いる。

すると、

\begin{align*}
S(L)
&=
k_B
\Biggl[
N\log N
-\frac{N+L}{2}
\log\frac{N+L}{2}
-\frac{N-L}{2}
\log\frac{N-L}{2}
\Biggr]
\end{align*}

となる。

ここで設問の指示に従い、

\begin{align*}
\frac{dS}{dL}
\simeq
\frac12
\left[
S(L)-S(L-2)
\right]
\end{align*}

を用いる。 【1-e08515】

上式から計算すると、

\begin{align*}
\frac{dS}{dL}
=
\frac{k_B}{2}
\log
\left(
\frac{N-L}{N+L}
\right)
\end{align*}

を得る。

一方、

\begin{align*}
U=-LmH
\end{align*}

より、

\begin{align*}
\frac{dU}{dL}
=
-mH
\end{align*}

である。

熱力学の定義

\begin{align*}
\frac{1}{T}
=
\left(
\frac{\partial S}{\partial U}
\right)
\end{align*}

を用いると、

\begin{align*}
\frac1T
&=
\frac{dS/dL}{dU/dL}
\\
&=
-\frac{k_B}{2mH}
\log
\left(
\frac{N-L}{N+L}
\right)
\\
&=
\frac{k_B}{2mH}
\log
\left(
\frac{N+L}{N-L}
\right)。
\end{align*}

したがって、

\begin{align*}
\frac{2mH}{k_BT}
=
\log
\left(
\frac{N+L}{N-L}
\right)
\end{align*}

を得る。

指数関数を取ると、

\begin{align*}
e^{2mH/(k_BT)}
=
\frac{N+L}{N-L}
\end{align*}

である。

これを $L$ について解けば、

\begin{align*}
L
=
N
\frac{
e^{2mH/(k_BT)}-1
}{
e^{2mH/(k_BT)}+1
}
=
N\tanh
\left(
\frac{mH}{k_BT}
\right)
\end{align*}

となる。

よって内部エネルギーは

\begin{align*}
U
=
-mHL
=
-mHN
\tanh
\left(
\frac{mH}{k_BT}
\right)
\end{align*}

と書ける。

以上より、

\begin{align*}
U(T,H)
=
-mHN
\tanh
\left(
\frac{mH}{k_BT}
\right)
\end{align*}

が示された。 【1-e08515】

\subsection{(iii)}

磁化 $M$ を求める。

熱力学の定義より、

\begin{align*}
M
=
-\left(
\frac{\partial U}{\partial H}
\right)_T
\end{align*}

である。

小問(ii)で得た結果から、

\begin{align*}
M
&=
mN
\tanh
\left(
\frac{mH}{k_BT}
\right)。
\end{align*}

ここで

\begin{align*}
\frac{mH}{k_BT}
\ll
1
\end{align*}

の極限を考える。

双曲線正接関数を展開すると、

\begin{align*}
\tanh x
=
x+O(x^3)
\end{align*}

であるから、

\begin{align*}
M
&\simeq
mN
\frac{mH}{k_BT}
\\
&=
\frac{Nm^2}{k_B}
\frac{H}{T}。
\end{align*}

したがって、

\begin{align*}
M
=
\frac{CH}{T}
\end{align*}

ただし

\begin{align*}
C=\frac{Nm^2}{k_B}
\end{align*}

である。

以上より、

\begin{align*}
mH
\ll
k_BT
\end{align*}

の高温・弱磁場極限で Curie の法則

\begin{align*}
M=\frac{CH}{T}
\end{align*}

が導かれた。
\section{5-5}

空洞放射について、内部エネルギー $U$、エントロピー $S$、体積 $V$、放射圧 $P$ の関係を調べる。 【1-d2e143】

\subsection{(i)}

示すべき式は

\begin{align*}
\left(
\frac{\partial U}{\partial V}
\right)_T
=
T
\left(
\frac{\partial S}{\partial V}
\right)_T
-P
\end{align*}

である。 【1-d2e143】

熱力学の基本関係式

\begin{align*}
dU
=
TdS-PdV
\end{align*}

を用いる。

エントロピーを $T,V$ の関数とみなすと、

\begin{align*}
dS
=
\left(
\frac{\partial S}{\partial T}
\right)_V dT
+
\left(
\frac{\partial S}{\partial V}
\right)_T dV
\end{align*}

である。

これを基本関係式へ代入すると、

\begin{align*}
dU
&=
T
\left[
\left(
\frac{\partial S}{\partial T}
\right)_V dT
+
\left(
\frac{\partial S}{\partial V}
\right)_T dV
\right]
-PdV \\
&=
T
\left(
\frac{\partial S}{\partial T}
\right)_V dT
+
\left[
T
\left(
\frac{\partial S}{\partial V}
\right)_T
-P
\right]dV。
\end{align*}

一方、

\begin{align*}
dU
=
\left(
\frac{\partial U}{\partial T}
\right)_V dT
+
\left(
\frac{\partial U}{\partial V}
\right)_T dV
\end{align*}

でもある。

したがって、$dV$ の係数を比較すると、

\begin{align*}
\left(
\frac{\partial U}{\partial V}
\right)_T
=
T
\left(
\frac{\partial S}{\partial V}
\right)_T
-P
\end{align*}

を得る。

以上より示すべき式が示された。 【1-d2e143】

\subsection{(ii)}

古典電磁気学より、

\begin{align*}
P=\frac{u}{3}
\end{align*}

が成り立つ。ここで $u=u(T)$ はエネルギー密度である。 【1-d2e143】

内部エネルギーは

\begin{align*}
U=u(T)V
\end{align*}

と書ける。

したがって、

\begin{align*}
\left(
\frac{\partial U}{\partial V}
\right)_T
=
u(T)
\end{align*}

である。

また、

\begin{align*}
F=U-TS
\end{align*}

について

\begin{align*}
dF=-SdT-PdV
\end{align*}

であるから Maxwell の関係式

\begin{align*}
\left(
\frac{\partial S}{\partial V}
\right)_T
=
\left(
\frac{\partial P}{\partial T}
\right)_V
\end{align*}

が成り立つ。

ここで

\begin{align*}
P=\frac13u(T)
\end{align*}

なので、

\begin{align*}
\left(
\frac{\partial P}{\partial T}
\right)_V
=
\frac13
\frac{du}{dT}
\end{align*}

である。

したがって、

\begin{align*}
\left(
\frac{\partial S}{\partial V}
\right)_T
=
\frac13
\frac{du}{dT}
\end{align*}

を得る。

これを小問(i)の式へ代入すると、

\begin{align*}
u
=
\frac{T}{3}
\frac{du}{dT}
-\frac{u}{3}
\end{align*}

となる。

整理すると、

\begin{align*}
4u
=
T
\frac{du}{dT}
\end{align*}

を得る。

したがって、

\begin{align*}
\frac{du}{u}
=
4
\frac{dT}{T}
\end{align*}

である。

積分すると、

\begin{align*}
\log u
=
4\log T+\mathrm{const.}
\end{align*}

となり、

\begin{align*}
u=aT^4
\end{align*}

を得る。ただし $a$ は定数である。

したがって、

\begin{align*}
U=uV
=
aVT^4
\end{align*}

となる。

以上より、

\begin{align*}
U=aVT^4
\end{align*}

が示された。これは Stefan-Boltzmann 則である。 【1-d2e143】

\subsection{(iii)}

エントロピーを求める。

基本関係式

\begin{align*}
dU
=
TdS-PdV
\end{align*}

を用いる。

まず、

\begin{align*}
U=aVT^4
\end{align*}

より、

\begin{align*}
dU
=
4aVT^3\,dT
+
aT^4\,dV
\end{align*}

である。

また、

\begin{align*}
P=\frac13u
=
\frac13aT^4
\end{align*}

である。

したがって、

\begin{align*}
TdS
&=
dU+PdV \\
&=
4aVT^3\,dT
+
aT^4\,dV
+
\frac13aT^4\,dV \\
&=
4aVT^3\,dT
+
\frac43aT^4\,dV。
\end{align*}

両辺を $T$ で割ると、

\begin{align*}
dS
=
4aVT^2\,dT
+
\frac43aT^3\,dV。
\end{align*}

ここで

\begin{align*}
d
\left(
\frac43aVT^3
\right)
&=
\frac43a
\left(
3VT^2\,dT
+
T^3\,dV
\right) \\
&=
4aVT^2\,dT
+
\frac43aT^3\,dV
\end{align*}

である。

したがって、

\begin{align*}
dS
=
d
\left(
\frac43aVT^3
\right)
\end{align*}

を得る。

積分すると、

\begin{align*}
S
=
\frac43aVT^3+\mathrm{const.}
\end{align*}

となる。

条件

\begin{align*}
S=0
\qquad
(T=0)
\end{align*}

を用いると積分定数は零である。

よって、

\begin{align*}
S
=
\frac43aVT^3
\end{align*}

を得る。

以上より、空洞放射のエントロピーは

\begin{align*}
S
=
\frac43aVT^3
\end{align*}

で与えられることが示された。
\section{5-6}

長さ $L$ のゴム紐について、

\begin{align*}
dU
=
TdS+\sigma dL
\end{align*}

が成り立つものとする。

また、理想ゴムでは張力が

\begin{align*}
\sigma=A(L)T,
\qquad
A(L)>0
\end{align*}

と書けるものとする。 【1-3bf4d0】

\subsection{(i)}

定長比熱

\begin{align*}
C_L
=
\left(
\frac{\partial U}{\partial T}
\right)_L
\end{align*}

が温度のみの関数であることを示す。 【1-3bf4d0】

まず内部エネルギーを $T,L$ の関数とみなす。

すると

\begin{align*}
dU
=
\left(
\frac{\partial U}{\partial T}
\right)_L dT
+
\left(
\frac{\partial U}{\partial L}
\right)_T dL
\end{align*}

である。

一方、

\begin{align*}
dS
=
\left(
\frac{\partial S}{\partial T}
\right)_L dT
+
\left(
\frac{\partial S}{\partial L}
\right)_T dL
\end{align*}

であるから、

\begin{align*}
dU
&=
T
\left[
\left(
\frac{\partial S}{\partial T}
\right)_L dT
+
\left(
\frac{\partial S}{\partial L}
\right)_T dL
\right]
+\sigma dL \\
&=
T
\left(
\frac{\partial S}{\partial T}
\right)_L dT
+
\left[
T
\left(
\frac{\partial S}{\partial L}
\right)_T
+\sigma
\right]dL。
\end{align*}

よって

\begin{align*}
\left(
\frac{\partial U}{\partial L}
\right)_T
=
T
\left(
\frac{\partial S}{\partial L}
\right)_T
+\sigma
\end{align*}

を得る。

次に Helmholtz 自由エネルギー

\begin{align*}
F=U-TS
\end{align*}

を導入する。

全微分は

\begin{align*}
dF
=
-SdT+\sigma dL
\end{align*}

である。

したがって Maxwell 関係式

\begin{align*}
\left(
\frac{\partial S}{\partial L}
\right)_T
=
-\left(
\frac{\partial \sigma}{\partial T}
\right)_L
\end{align*}

を得る。

理想ゴムでは

\begin{align*}
\sigma=A(L)T
\end{align*}

なので、

\begin{align*}
\left(
\frac{\partial \sigma}{\partial T}
\right)_L
=
A(L)
\end{align*}

である。

したがって

\begin{align*}
\left(
\frac{\partial S}{\partial L}
\right)_T
=
-A(L)
\end{align*}

となる。

これを代入すると、

\begin{align*}
\left(
\frac{\partial U}{\partial L}
\right)_T
&=
-TA(L)+A(L)T \\
&=
0。
\end{align*}

ゆえに、

\begin{align*}
U=U(T)
\end{align*}

である。

したがって、

\begin{align*}
C_L
=
\left(
\frac{\partial U}{\partial T}
\right)_L
=
\frac{dU(T)}{dT}
\end{align*}

となり、$C_L$ は温度のみの関数である。

以上より、

\begin{align*}
C_L=C_L(T)
\end{align*}

が示された。 【1-3bf4d0】

\subsection{(ii)}

温度一定でゴム紐を引き延ばす。

前問で得た Maxwell 関係式

\begin{align*}
\left(
\frac{\partial S}{\partial L}
\right)_T
=
-\left(
\frac{\partial \sigma}{\partial T}
\right)_L
\end{align*}

を用いる。

理想ゴムでは

\begin{align*}
\sigma=A(L)T
\end{align*}

であるから、

\begin{align*}
\left(
\frac{\partial \sigma}{\partial T}
\right)_L
=
A(L)
>
0
\end{align*}

である。

したがって、

\begin{align*}
\left(
\frac{\partial S}{\partial L}
\right)_T
=
-A(L)
<
0
\end{align*}

を得る。

よって長さを増加させるとエントロピーは減少する。

以上より、温度一定で理想ゴムを引き延ばすとエントロピーは減少する。 【1-3bf4d0】

\subsection{(iii)}

準静的かつ断熱的にゴム紐を引き延ばす場合を考える。

断熱条件より、

\begin{align*}
dS=0
\end{align*}

である。

エントロピーの全微分

\begin{align*}
dS
=
\left(
\frac{\partial S}{\partial T}
\right)_L dT
+
\left(
\frac{\partial S}{\partial L}
\right)_T dL
\end{align*}

を用いる。

ここで

\begin{align*}
C_L
=
T
\left(
\frac{\partial S}{\partial T}
\right)_L
\end{align*}

より、

\begin{align*}
\left(
\frac{\partial S}{\partial T}
\right)_L
=
\frac{C_L}{T}
\end{align*}

である。

また前問より、

\begin{align*}
\left(
\frac{\partial S}{\partial L}
\right)_T
=
-A(L)
\end{align*}

である。

したがって、

\begin{align*}
0
=
\frac{C_L}{T}dT
-
A(L)dL
\end{align*}

を得る。

整理すると、

\begin{align*}
dT
=
\frac{TA(L)}{C_L}
\,dL。
\end{align*}

ここで

\begin{align*}
T>0,
\qquad
A(L)>0,
\qquad
C_L>0
\end{align*}

である。

よって、

\begin{align*}
dL>0
\end{align*}

ならば

\begin{align*}
dT>0
\end{align*}

である。

以上より、理想ゴムを断熱的に引き延ばすと温度は上昇する。 【1-3bf4d0】

\subsection{(iv)}

張力一定の条件で温度を上昇させる。

張力は

\begin{align*}
\sigma
=
A(L)T
\end{align*}

である。

全微分を取ると、

\begin{align*}
d\sigma
=
A'(L)T\,dL
+
A(L)dT
\end{align*}

である。

張力一定では

\begin{align*}
d\sigma=0
\end{align*}

なので、

\begin{align*}
A'(L)T\,dL
+
A(L)dT
=
0
\end{align*}

となる。

したがって、

\begin{align*}
dL
=
-\frac{A(L)}{A'(L)T}
dT。
\end{align*}

問題の仮定より、

\begin{align*}
A'(L)>0,
\qquad
A(L)>0,
\qquad
T>0
\end{align*}

である。

ゆえに、

\begin{align*}
dT>0
\end{align*}

のとき

\begin{align*}
dL<0
\end{align*}

となる。

以上より、張力を一定に保ちながら理想ゴムの温度を上昇させると、ゴム紐は縮む。

これは加熱によってゴムの配向の乱雑さが増し、エントロピー増大の方向へ向かうためである。
\section{6-1}

熱平衡状態において

\begin{align*}
\left(
\frac{\partial^2 U}{\partial S^2}
\right)_V
&>0,
\\
\left(
\frac{\partial^2 U}{\partial V^2}
\right)_S
&>0,
\\
\left(
\frac{\partial^2 U}{\partial S^2}
\right)_V
\left(
\frac{\partial^2 U}{\partial V^2}
\right)_S
&>
\left(
\frac{\partial^2 U}{\partial S\partial V}
\right)^2
\end{align*}

が成り立つと仮定する。これらを熱力学不等式と呼ぶ。 【1-8f09c1】

\subsection{(i)}

断熱的に体積を増加させると圧力が低下することを示す。

熱力学の基本関係式

\begin{align*}
dU
=
TdS-pdV
\end{align*}

より、

\begin{align*}
T
=
\left(
\frac{\partial U}{\partial S}
\right)_V,
\qquad
p
=
-\left(
\frac{\partial U}{\partial V}
\right)_S
\end{align*}

である。

したがって、

\begin{align*}
\left(
\frac{\partial p}{\partial V}
\right)_S
=
-
\left(
\frac{\partial^2U}{\partial V^2}
\right)_S。
\end{align*}

熱力学不等式より

\begin{align*}
\left(
\frac{\partial^2U}{\partial V^2}
\right)_S
>0
\end{align*}

であるから、

\begin{align*}
\left(
\frac{\partial p}{\partial V}
\right)_S
<0
\end{align*}

を得る。

よって断熱的に体積が増加すると圧力は低下する。

以上より、

\begin{align*}
\left(
\frac{\partial p}{\partial V}
\right)_S
<0
\end{align*}

が示された。 【1-8f09c1】

\subsection{(ii)}

定積比熱が正であることを示す。

定積比熱は

\begin{align*}
C_V
=
T
\left(
\frac{\partial S}{\partial T}
\right)_V
\end{align*}

で定義される。

ここで

\begin{align*}
T
=
\left(
\frac{\partial U}{\partial S}
\right)_V
\end{align*}

を $S$ で微分すると、

\begin{align*}
\left(
\frac{\partial T}{\partial S}
\right)_V
=
\left(
\frac{\partial^2U}{\partial S^2}
\right)_V
\end{align*}

となる。

熱力学不等式より、

\begin{align*}
\left(
\frac{\partial T}{\partial S}
\right)_V
>0
\end{align*}

である。

したがって、

\begin{align*}
\left(
\frac{\partial S}{\partial T}
\right)_V
=
\left[
\left(
\frac{\partial T}{\partial S}
\right)_V
\right]^{-1}
>0
\end{align*}

を得る。

さらに平衡熱力学では

\begin{align*}
T>0
\end{align*}

であるから、

\begin{align*}
C_V
=
T
\left(
\frac{\partial S}{\partial T}
\right)_V
>0
\end{align*}

となる。

以上より、

\begin{align*}
C_V>0
\end{align*}

が示された。 【1-8f09c1】

\subsection{(iii)}

定圧比熱が定積比熱より大きいことを示す。

まず Jacobian の恒等式

\begin{align*}
\left(
\frac{\partial p}{\partial V}
\right)_T
=
\frac{
\partial(p,T)/\partial(V,S)
}{
\partial(T,V)/\partial(V,S)
}
\end{align*}

を用いる。

分子を計算すると、

\begin{align*}
\frac{\partial(p,T)}{\partial(V,S)}
&=
\left(
\frac{\partial p}{\partial V}
\right)_S
\left(
\frac{\partial T}{\partial S}
\right)_V
-
\left(
\frac{\partial p}{\partial S}
\right)_V
\left(
\frac{\partial T}{\partial V}
\right)_S。
\end{align*}

ここで

\begin{align*}
p
=
-\left(
\frac{\partial U}{\partial V}
\right)_S,
\qquad
T
=
\left(
\frac{\partial U}{\partial S}
\right)_V
\end{align*}

を用いると、

\begin{align*}
\left(
\frac{\partial p}{\partial S}
\right)_V
=
-
\left(
\frac{\partial^2U}{\partial S\partial V}
\right),
\qquad
\left(
\frac{\partial T}{\partial V}
\right)_S
=
\left(
\frac{\partial^2U}{\partial V\partial S}
\right)
\end{align*}

である。

したがって、

\begin{align*}
\frac{\partial(p,T)}{\partial(V,S)}
&=
-
\left[
\left(
\frac{\partial^2U}{\partial S^2}
\right)_V
\left(
\frac{\partial^2U}{\partial V^2}
\right)_S
-
\left(
\frac{\partial^2U}{\partial S\partial V}
\right)^2
\right]。
\end{align*}

熱力学不等式より角括弧内は正であるから、

\begin{align*}
\frac{\partial(p,T)}{\partial(V,S)}
<0
\end{align*}

となる。

また、

\begin{align*}
\frac{\partial(T,V)}{\partial(V,S)}
=
-
\left(
\frac{\partial T}{\partial S}
\right)_V
<0
\end{align*}

である。

よって、

\begin{align*}
\left(
\frac{\partial p}{\partial V}
\right)_T
<0
\end{align*}

を得る。

したがって等温圧縮率

\begin{align*}
\kappa_T
=
-\frac1V
\left(
\frac{\partial V}{\partial p}
\right)_T
\end{align*}

は正である。

第2回演習で導いた関係式

\begin{align*}
\frac{C_P}{C_V}
=
\frac{\kappa_T}{\kappa_S}
\end{align*}

を用いる。 【1-8f09c1】

小問(i)より、

\begin{align*}
\left(
\frac{\partial p}{\partial V}
\right)_S
<0
\end{align*}

であるので、

\begin{align*}
\kappa_S>0
\end{align*}

である。

したがって、

\begin{align*}
\frac{C_P}{C_V}
=
\frac{\kappa_T}{\kappa_S}
>0。
\end{align*}

さらに熱力学不等式から

\begin{align*}
\kappa_T>\kappa_S
\end{align*}

が従うので、

\begin{align*}
\frac{C_P}{C_V}
>1
\end{align*}

となる。

小問(ii)で

\begin{align*}
C_V>0
\end{align*}

を示したから、

\begin{align*}
C_P>C_V
\end{align*}

を得る。

以上より、定圧比熱は定積比熱より常に大きいことが示された。
\section{6-2}

等温圧縮率

\begin{align*}
\kappa_T
=
-\frac{1}{V}
\left(
\frac{\partial V}{\partial p}
\right)_T
\end{align*}

および断熱圧縮率

\begin{align*}
\kappa_S
=
-\frac{1}{V}
\left(
\frac{\partial V}{\partial p}
\right)_S
\end{align*}

を考える。 【1-57ef3f】

\subsection{(i)}

まず

\begin{align*}
\kappa_T-\kappa_S
=
-\frac1V
\left[
\left(
\frac{\partial V}{\partial p}
\right)_T
-
\left(
\frac{\partial V}{\partial p}
\right)_S
\right]
\end{align*}

である。

ここで一般の偏微分公式

\begin{align*}
\left(
\frac{\partial V}{\partial p}
\right)_T
=
\left(
\frac{\partial V}{\partial p}
\right)_S
+
\left(
\frac{\partial V}{\partial S}
\right)_p
\left(
\frac{\partial S}{\partial p}
\right)_T
\end{align*}

を用いると、

\begin{align*}
\kappa_T-\kappa_S
=
-\frac1V
\left(
\frac{\partial V}{\partial S}
\right)_p
\left(
\frac{\partial S}{\partial p}
\right)_T
\end{align*}

を得る。

次に、

\begin{align*}
dV
=
\left(
\frac{\partial V}{\partial p}
\right)_S dp
+
\left(
\frac{\partial V}{\partial S}
\right)_p dS
\end{align*}

を考える。

体積一定、

\begin{align*}
dV=0
\end{align*}

とすると、

\begin{align*}
\left(
\frac{\partial p}{\partial S}
\right)_V
=
-
\frac{
\left(
\frac{\partial V}{\partial S}
\right)_p
}{
\left(
\frac{\partial V}{\partial p}
\right)_S
}
\end{align*}

である。

一方、Jacobian の恒等式

\begin{align*}
\left(
\frac{\partial V}{\partial S}
\right)_p
\left(
\frac{\partial S}{\partial p}
\right)_V
\left(
\frac{\partial p}{\partial V}
\right)_S
=
-1
\end{align*}

より、

\begin{align*}
\left(
\frac{\partial V}{\partial S}
\right)_p
=
-
\left(
\frac{\partial T}{\partial p}
\right)_S
\end{align*}

が成り立つ。

したがって、

\begin{align*}
\kappa_T-\kappa_S
=
-\frac1V
\left(
\frac{\partial T}{\partial p}
\right)_S
\left(
\frac{\partial S}{\partial p}
\right)_T
\end{align*}

を得る。

以上より、

\begin{align*}
\boxed{
\kappa_T-\kappa_S
=
-\frac1V
\left(
\frac{\partial T}{\partial p}
\right)_S
\left(
\frac{\partial S}{\partial p}
\right)_T
}
\end{align*}

が示された。 【1-57ef3f】

\subsection{(ii)}

前問で得た式

\begin{align*}
\kappa_T-\kappa_S
=
-\frac1V
\left(
\frac{\partial T}{\partial p}
\right)_S
\left(
\frac{\partial S}{\partial p}
\right)_T
\end{align*}

を用いる。

まず Maxwell の関係式

\begin{align*}
\left(
\frac{\partial S}{\partial p}
\right)_T
=
-
\left(
\frac{\partial V}{\partial T}
\right)_p
\end{align*}

より、

\begin{align*}
\kappa_T-\kappa_S
=
\frac1V
\left(
\frac{\partial T}{\partial p}
\right)_S
\left(
\frac{\partial V}{\partial T}
\right)_p
\end{align*}

となる。

ここで循環関係

\begin{align*}
\left(
\frac{\partial T}{\partial p}
\right)_S
=
-
\left(
\frac{\partial T}{\partial S}
\right)_p
\left(
\frac{\partial S}{\partial p}
\right)_T
\end{align*}

を用い、さらに

\begin{align*}
\left(
\frac{\partial T}{\partial S}
\right)_p
=
\left[
\left(
\frac{\partial S}{\partial T}
\right)_p
\right]^{-1}
=
\frac{T}{C_P}
\end{align*}

を代入すると、

\begin{align*}
\kappa_T-\kappa_S
=
\frac{T}{VC_P}
\left[
\left(
\frac{\partial V}{\partial T}
\right)_p
\right]^2
\end{align*}

を得る。

熱平衡状態では

\begin{align*}
T>0,
\qquad
V>0,
\qquad
C_P>0
\end{align*}

である。

また、

\begin{align*}
\left[
\left(
\frac{\partial V}{\partial T}
\right)_p
\right]^2
\ge 0
\end{align*}

であるから、

\begin{align*}
\kappa_T-\kappa_S
\ge 0
\end{align*}

を得る。

通常の物質では

\begin{align*}
\left(
\frac{\partial V}{\partial T}
\right)_p
\neq 0
\end{align*}

であるので、

\begin{align*}
\kappa_T-\kappa_S
>
0
\end{align*}

となる。

したがって、

\begin{align*}
\boxed{
\kappa_T>\kappa_S
}
\end{align*}

が示された。

これは定温で圧縮するときの方が、断熱圧縮よりも気体が圧縮されやすいことを意味している。
\section{6-3}

Joule-Thomson 実験を考える。断熱材で覆われた容器の内部を多孔質壁で二つに分け、気体を高圧側から低圧側へ流す。このとき、Joule-Thomson 係数

\begin{align*}
\xi
=
\left(
\frac{\partial T}{\partial p}
\right)_H
\end{align*}

を求める。 【1-622fa9】

\subsection{(i)}

まず、この過程でエンタルピーが保存されることを示す。

多孔質壁を通過する気体について定常状態を考える。

高圧側から低圧側へ体積

\begin{align*}
\Delta V_1
\end{align*}

の気体を押し込むために外界がする仕事は

\begin{align*}
W_1=p_1\Delta V_1
\end{align*}

である。

一方、低圧側では体積

\begin{align*}
\Delta V_2
\end{align*}

だけ膨張するので、気体が外界へする仕事は

\begin{align*}
W_2=p_2\Delta V_2
\end{align*}

である。

したがって気体に対する正味の仕事は

\begin{align*}
W=p_1\Delta V_1-p_2\Delta V_2
\end{align*}

となる。

この過程は断熱的なので、

\begin{align*}
Q=0
\end{align*}

である。

熱力学第1法則

\begin{align*}
\Delta U=W+Q
\end{align*}

より、

\begin{align*}
U_2-U_1
=
p_1\Delta V_1-p_2\Delta V_2
\end{align*}

を得る。

整理すると、

\begin{align*}
U_1+p_1\Delta V_1
=
U_2+p_2\Delta V_2。
\end{align*}

ここでエンタルピー

\begin{align*}
H=U+pV
\end{align*}

を用いると、

\begin{align*}
H_1=H_2
\end{align*}

となる。

以上より、Joule-Thomson 実験ではエンタルピーが保存される。 【1-622fa9】

\subsection{(ii)}

Joule-Thomson 係数

\begin{align*}
\xi
=
\left(
\frac{\partial T}{\partial p}
\right)_H
\end{align*}

を求める。

エンタルピーを変数 $T,p$ の関数とみなして全微分すると、

\begin{align*}
dH
=
\left(
\frac{\partial H}{\partial T}
\right)_p dT
+
\left(
\frac{\partial H}{\partial p}
\right)_T dp。
\end{align*}

ここで

\begin{align*}
\left(
\frac{\partial H}{\partial T}
\right)_p
=
C_P
\end{align*}

であるから、

\begin{align*}
dH
=
C_P dT
+
\left(
\frac{\partial H}{\partial p}
\right)_T dp
\end{align*}

となる。

Joule-Thomson 過程では

\begin{align*}
dH=0
\end{align*}

なので、

\begin{align*}
\left(
\frac{\partial T}{\partial p}
\right)_H
=
-\frac1{C_P}
\left(
\frac{\partial H}{\partial p}
\right)_T
\end{align*}

を得る。

したがって、

\begin{align*}
\xi
=
-\frac1{C_P}
\left(
\frac{\partial H}{\partial p}
\right)_T
\end{align*}

である。

次に、

\begin{align*}
H=U+pV
\end{align*}

を微分する。

\begin{align*}
dH
=
dU+p\,dV+V\,dp。
\end{align*}

また、

\begin{align*}
dU
=
TdS-p\,dV
\end{align*}

であるから、

\begin{align*}
dH
=
TdS+Vdp。
\end{align*}

さらにエントロピーを $T,p$ の関数とみなすと、

\begin{align*}
dS
=
\left(
\frac{\partial S}{\partial T}
\right)_p dT
+
\left(
\frac{\partial S}{\partial p}
\right)_T dp
\end{align*}

である。

これを代入すると、

\begin{align*}
dH
&=
T
\left(
\frac{\partial S}{\partial T}
\right)_p dT
+
\left[
T
\left(
\frac{\partial S}{\partial p}
\right)_T
+
V
\right]dp。
\end{align*}

したがって、

\begin{align*}
\left(
\frac{\partial H}{\partial p}
\right)_T
=
T
\left(
\frac{\partial S}{\partial p}
\right)_T
+
V
\end{align*}

を得る。

ここで Gibbs 自由エネルギー

\begin{align*}
G=H-TS
\end{align*}

の全微分

\begin{align*}
dG
=
-SdT+Vdp
\end{align*}

から得られる Maxwell 関係式

\begin{align*}
\left(
\frac{\partial S}{\partial p}
\right)_T
=
-
\left(
\frac{\partial V}{\partial T}
\right)_p
\end{align*}

を用いると、

\begin{align*}
\left(
\frac{\partial H}{\partial p}
\right)_T
=
V-T
\left(
\frac{\partial V}{\partial T}
\right)_p
\end{align*}

となる。

よって、

\begin{align*}
\xi
&=
-\frac1{C_P}
\left[
V-T
\left(
\frac{\partial V}{\partial T}
\right)_p
\right]
\\
&=
\frac{
T
\left(
\frac{\partial V}{\partial T}
\right)_p
-V
}{
C_P
}。
\end{align*}

したがって、

\begin{align*}
\boxed{
\xi
=
\left(
\frac{\partial T}{\partial p}
\right)_H
=
\frac{
T
\left(
\frac{\partial V}{\partial T}
\right)_p
-V
}{
C_P
}
}
\end{align*}

が示された。 【1-622fa9】

\subsection{(iii)}

理想気体では

\begin{align*}
pV=NRT
\end{align*}

が成り立つ。

したがって、

\begin{align*}
V
=
\frac{NRT}{p}
\end{align*}

であり、

\begin{align*}
\left(
\frac{\partial V}{\partial T}
\right)_p
=
\frac{NR}{p}
=
\frac{V}{T}
\end{align*}

を得る。

これを Joule-Thomson 係数の式へ代入すると、

\begin{align*}
\xi
&=
\frac{
T(V/T)-V
}{
C_P
}
\\
&=
0。
\end{align*}

したがって、

\begin{align*}
\boxed{
\xi=0
}
\end{align*}

となる。

以上より、理想気体では多孔質壁を通過しても温度は変化しない。

これは理想気体では内部エネルギーが温度のみに依存し、分子間相互作用によるエネルギー交換が存在しないためである。
\section{6-4}

van der Waals 気体について Joule-Thomson 係数および逆転温度を求める。状態方程式は

\begin{align*}
\left(
p+\frac{aN^2}{V^2}
\right)
(V-bN)
=
NRT
\end{align*}

で与えられるものとする。 【1-428a98】

\subsection{(i)}

まず Joule-Thomson 係数

\begin{align*}
\xi
=
\left(
\frac{\partial T}{\partial p}
\right)_H
\end{align*}

を求める。

前問より、

\begin{align*}
\xi
=
\frac{
T
\left(
\frac{\partial V}{\partial T}
\right)_p
-V
}{
C_P
}
\end{align*}

である。 【1-428a98】

したがって、

\begin{align*}
\left(
\frac{\partial V}{\partial T}
\right)_p
\end{align*}

を求めればよい。

状態方程式

\begin{align*}
f(T,p,V)
=
\left(
p+\frac{aN^2}{V^2}
\right)
(V-bN)
-NRT
=
0
\end{align*}

を温度で偏微分すると、

\begin{align*}
\frac{\partial f}{\partial T}
+
\frac{\partial f}{\partial V}
\left(
\frac{\partial V}{\partial T}
\right)_p
=
0
\end{align*}

となる。

まず、

\begin{align*}
\frac{\partial f}{\partial T}
=
-NR
\end{align*}

である。

また、

\begin{align*}
\frac{\partial f}{\partial V}
&=
\left(
p+\frac{aN^2}{V^2}
\right)
-\frac{2aN^2}{V^3}(V-bN) \\
&=
p-\frac{aN^2}{V^2}
+\frac{2abN^3}{V^3}。
\end{align*}

したがって、

\begin{align*}
\left(
\frac{\partial V}{\partial T}
\right)_p
=
\frac{NR}
{
p-\dfrac{aN^2}{V^2}
+\dfrac{2abN^3}{V^3}
}。
\end{align*}

これをそのまま用いてもよいが、Joule-Thomson 効果は一般に低密度領域で測定されるので、

\begin{align*}
\frac{bN}{V}
\ll 1
\end{align*}

を仮定して一次まで展開する。

状態方程式から、

\begin{align*}
p
=
\frac{NRT}{V-bN}
-\frac{aN^2}{V^2}
\simeq
\frac{NRT}{V}
+
\frac{bN^2RT}{V^2}
-\frac{aN^2}{V^2}
\end{align*}

である。

このとき、

\begin{align*}
V
\simeq
\frac{NRT}{p}
+
\frac{Nb}{}
-\frac{aN}{RT}
\end{align*}

を得る。

したがって、

\begin{align*}
\left(
\frac{\partial V}{\partial T}
\right)_p
\simeq
\frac{NR}{p}
+
\frac{aN}{RT^2}
\end{align*}

である。

よって、

\begin{align*}
T
\left(
\frac{\partial V}{\partial T}
\right)_p
-V
&=
T
\left(
\frac{NR}{p}
+
\frac{aN}{RT^2}
\right)
-
\left(
\frac{NRT}{p}
+
Nb
-
\frac{aN}{RT}
\right)
\\
&=
\frac{2aN}{RT}
-Nb。
\end{align*}

したがって、

\begin{align*}
\boxed{
\xi
=
\frac{N}{C_P}
\left(
\frac{2a}{RT}
-b
\right)
}
\end{align*}

を得る。 【1-428a98】

\subsection{(ii)}

逆転温度 $T_{\mathrm{inv}}$ は

\begin{align*}
\xi=0
\end{align*}

を満たす温度である。

前問の結果より、

\begin{align*}
\frac{N}{C_P}
\left(
\frac{2a}{RT_{\mathrm{inv}}}
-b
\right)
=
0
\end{align*}

である。

\begin{align*}
N>0,
\qquad
C_P>0
\end{align*}

であるから、

\begin{align*}
\frac{2a}{RT_{\mathrm{inv}}}
-b
=
0
\end{align*}

を得る。

したがって、

\begin{align*}
T_{\mathrm{inv}}
=
\frac{2a}{Rb}
\end{align*}

となる。

よって、

\begin{align*}
T
<
\frac{2a}{Rb}
\end{align*}

では

\begin{align*}
\xi>0
\end{align*}

となり、多孔質壁通過後に温度は上昇する。

一方、

\begin{align*}
T
>
\frac{2a}{Rb}
\end{align*}

では

\begin{align*}
\xi<0
\end{align*}

となり、温度は低下する。

以上より、van der Waals 気体の逆転温度は

\begin{align*}
T_{\mathrm{inv}}
=
\frac{2a}{Rb}
\end{align*}

で与えられる。
\section{6-5}

Gibbs自由エネルギー

\begin{align*}
G=G(T,p,N)
\end{align*}

を用いて気液相転移を考える。

\subsection{(i)}

Gibbs自由エネルギーの全微分は

\begin{align*}
dG
=
-S\,dT+V\,dp
\end{align*}

である。

したがって、

\begin{align*}
S
=
-\left(
\frac{\partial G}{\partial T}
\right)_p,
\qquad
V
=
\left(
\frac{\partial G}{\partial p}
\right)_T
\end{align*}

を得る。

一次相転移では相境界上で Gibbs 自由エネルギーそのものは連続であるが、その一次微分は不連続となる。

よって、

\begin{align*}
S_G\neq S_L,
\qquad
V_G\neq V_L
\end{align*}

である。

以上より、一次相転移ではエントロピーと体積が不連続に変化する。

\subsection{(ii)}

相境界上の点

\begin{align*}
(T,p)
\end{align*}

から

\begin{align*}
(T+dT,p+dp)
\end{align*}

へ移ることを考える。

液相について、

\begin{align*}
dG_L
=
-S_L\,dT+V_L\,dp
\end{align*}

である。

気相についても同様に、

\begin{align*}
dG_G
=
-S_G\,dT+V_G\,dp
\end{align*}

である。

相境界上では常に

\begin{align*}
G_L=G_G
\end{align*}

であるから、

\begin{align*}
dG_L=dG_G
\end{align*}

が成り立つ。

したがって、

\begin{align*}
-S_L\,dT+V_L\,dp
=
-S_G\,dT+V_G\,dp
\end{align*}

を得る。

整理すると、

\begin{align*}
(V_G-V_L)\,dp
=
(S_G-S_L)\,dT
\end{align*}

となる。

よって、

\begin{align*}
\frac{dp}{dT}
=
\frac{S_G-S_L}{V_G-V_L}
\end{align*}

を得る。

\subsection{(iii)}

気化潜熱を

\begin{align*}
L
=
T(S_G-S_L)
\end{align*}

と定義する。

したがって、

\begin{align*}
S_G-S_L
=
\frac{L}{T}
\end{align*}

である。

これを前問の結果へ代入すると、

\begin{align*}
\frac{dp}{dT}
=
\frac{L}
{T(V_G-V_L)}
\end{align*}

を得る。

以上より、

\begin{align*}
\frac{dp}{dT}
=
\frac{L}
{T(V_G-V_L)}
\end{align*}

が成り立つ。

これを Clapeyron-Clausius の関係式という。

\subsection{(iv)}

液体のモル体積は気体のモル体積に比べて十分小さいので、

\begin{align*}
V_G-V_L
\simeq
V_G
\end{align*}

と近似する。

さらに蒸気を理想気体とみなすと、

\begin{align*}
V_G
=
\frac{RT}{p}
\end{align*}

である。

したがって、

\begin{align*}
\frac{dp}{dT}
=
\frac{Lp}{RT^2}
\end{align*}

を得る。

変数分離すると、

\begin{align*}
\frac{dp}{p}
=
\frac{L}{R}
\frac{dT}{T^2}
\end{align*}

である。

両辺を積分すると、

\begin{align*}
\log p
=
-\frac{L}{RT}
+\mathrm{const.}
\end{align*}

となる。

よって、

\begin{align*}
p
=
p_0
\exp\left(
-\frac{L}{RT}
\right)
\end{align*}

を得る。

\subsection{(v)}

二つの沸点

\begin{align*}
(T_0,p_0),
\qquad
(T,p)
\end{align*}

について積分形

\begin{align*}
\log\frac{p}{p_0}
=
-\frac{L}{R}
\left(
\frac1T-\frac1{T_0}
\right)
\end{align*}

を用いる。

ここで

\begin{align*}
T_0
=
373.15\,\mathrm{K},
\qquad
T
=
368.15\,\mathrm{K},
\qquad
L
=
4.06\times10^4\,\mathrm{J/mol}
\end{align*}

を代入すると、

\begin{align*}
\log\frac{p}{1.013\times10^5}
=
-\frac{4.06\times10^4}{8.314}
\left(
\frac1{368.15}
-\frac1{373.15}
\right)
\end{align*}

となる。

計算すると、

\begin{align*}
p
\simeq
8.37\times10^4\,\mathrm{Pa}
\end{align*}

を得る。

したがって、

\begin{align*}
p
\simeq
837\,\mathrm{hPa}
\end{align*}

である。

標準気圧との差はおよそ

\begin{align*}
1013-837
=
176\,\mathrm{hPa}
\end{align*}

であるから、

\begin{align*}
h
\simeq
1.76\times10^3\,\mathrm{m}
\end{align*}

となる。

以上より、95 ℃で沸騰する場所の標高は約

\begin{align*}
1.8\times10^3\,\mathrm{m}
\end{align*}

と見積もられる。
\section{6-6}

van der Waals 気体

\begin{align*}
\left(
p+\frac{aN^2}{V^2}
\right)
(V-bN)
=
RNT
\end{align*}

について、相転移と臨界点を調べる。 【1-a4b144】

\subsection{(i)}

気液相転移に伴う潜熱を求める。

気液平衡では温度 $T$、圧力 $p$ が一定であり、相転移前後で Gibbs 自由エネルギーは等しい。

気化潜熱は

\begin{align*}
Q=T(S_G-S_L)
\end{align*}

で与えられる。 【1-a4b144】

第2回で導いた van der Waals 気体の断熱関係式から、

\begin{align*}
dS
=
cRN
\frac{dT}{T}
+
RN
\frac{dV}{V-bN}
\end{align*}

である。

相転移では温度一定なので

\begin{align*}
dT=0
\end{align*}

であり、

\begin{align*}
dS
=
RN
\frac{dV}{V-bN}
\end{align*}

となる。

液体相から気体相への変化について積分すると、

\begin{align*}
S_G-S_L
&=
RN
\int_{V_L}^{V_G}
\frac{dV}{V-bN}
\\
&=
RN
\log
\left(
\frac{V_G-bN}
     {V_L-bN}
\right)。
\end{align*}

したがって、

\begin{align*}
Q
&=
T(S_G-S_L)
\\
&=
RNT
\log
\left(
\frac{V_G-bN}
     {V_L-bN}
\right)
\end{align*}

を得る。

以上より、

\begin{align*}
Q
=
RNT
\log
\left(
\frac{V_G-bN}
     {V_L-bN}
\right)
\end{align*}

が示された。 【1-a4b144】

\subsection{(ii)}

臨界点では等温線が変曲点を持つ。

したがって、

\begin{align*}
\left(
\frac{\partial p}{\partial V}
\right)_{T=T_C}
=
0,
\qquad
\left(
\frac{\partial^2p}{\partial V^2}
\right)_{T=T_C}
=
0
\end{align*}

が成り立つ。 【1-a4b144】

まず状態方程式を

\begin{align*}
p
=
\frac{RNT}{V-bN}
-
\frac{aN^2}{V^2}
\end{align*}

と書く。

一階微分は

\begin{align*}
\frac{\partial p}{\partial V}
=
-\frac{RNT}{(V-bN)^2}
+
\frac{2aN^2}{V^3}
\end{align*}

である。

臨界点では

\begin{align*}
-\frac{RNT_C}{(V_C-bN)^2}
+
\frac{2aN^2}{V_C^3}
=
0。
\end{align*}

したがって、

\begin{align*}
RNT_C
=
\frac{2aN^2(V_C-bN)^2}{V_C^3}
\end{align*}

を得る。 \hfill (1)

次に二階微分を計算する：

\begin{align*}
\frac{\partial^2p}{\partial V^2}
=
\frac{2RNT}{(V-bN)^3}
-
\frac{6aN^2}{V^4}
\end{align*}

である。

臨界点条件より、

\begin{align*}
\frac{2RNT_C}{(V_C-bN)^3}
=
\frac{6aN^2}{V_C^4}
\end{align*}

すなわち

\begin{align*}
RNT_C
=
\frac{3aN^2(V_C-bN)^3}{V_C^4}
\end{align*}

を得る。 \hfill (2)

(1)、(2) を等しいとすると、

\begin{align*}
\frac{2(V_C-bN)^2}{V_C^3}
=
\frac{3(V_C-bN)^3}{V_C^4}
\end{align*}

となる。

整理すると、

\begin{align*}
2V_C
=
3(V_C-bN)
\end{align*}

である。

したがって、

\begin{align*}
V_C
=
3bN
\end{align*}

を得る。

これを (1) に代入すると、

\begin{align*}
RNT_C
&=
\frac{
2aN^2(2bN)^2
}{
(3bN)^3
}
\\
&=
\frac{8aN}{27b}
\end{align*}

となる。

よって、

\begin{align*}
T_C
=
\frac{8a}{27Rb}
\end{align*}

である。

最後に状態方程式へ代入すると、

\begin{align*}
p_C
&=
\frac{RNT_C}{V_C-bN}
-\frac{aN^2}{V_C^2}
\\
&=
\frac{
N
\left(
\dfrac{8a}{27b}
\right)
}{
2bN
}
-
\frac{aN^2}
{(3bN)^2}
\\
&=
\frac{4a}{27b^2}
-\frac{a}{9b^2}
\\
&=
\frac{a}{27b^2}。
\end{align*}

以上より、臨界点は

\begin{align*}
V_C
=
3bN,
\qquad
T_C
=
\frac{8a}{27Rb},
\qquad
p_C
=
\frac{a}{27b^2}
\end{align*}

で与えられる。

したがって van der Waals 気体の臨界定数は係数 $a,b$ のみによって決定されることが示された。 【1-a4b144】
\section{6-7}

理想常磁性体について、断熱消磁による冷却効果を調べる。常磁性体の基本関係式

\begin{align*}
dU
=
TdS+HdM
\end{align*}

および Curie の法則

\begin{align*}
M=\frac{CH}{T}
\end{align*}

が成り立つものとする。 【1-45ab9c】

\subsection{(i)}

まず

\begin{align*}
dS
=
\frac{C_H}{T}dT
+
\left(
\frac{\partial M}{\partial T}
\right)_H dH
\end{align*}

を示す。 【1-45ab9c】

Gibbs 自由エネルギー

\begin{align*}
G=U-TS-HM
\end{align*}

を導入する。

全微分を計算すると、

\begin{align*}
dG
&=
dU-TdS-SdT-HdM-MdH \\
&=
(TdS+HdM)-TdS-SdT-HdM-MdH \\
&=
-SdT-MdH。
\end{align*}

したがって、

\begin{align*}
S
=
-\left(
\frac{\partial G}{\partial T}
\right)_H,
\qquad
M
=
-\left(
\frac{\partial G}{\partial H}
\right)_T
\end{align*}

である。

混合偏導関数の交換可能性から、

\begin{align*}
\left(
\frac{\partial S}{\partial H}
\right)_T
=
\left(
\frac{\partial M}{\partial T}
\right)_H
\end{align*}

を得る。

これは Maxwell の関係式である。

一方、エントロピーを変数 $T,H$ の関数として考えると、

\begin{align*}
dS
=
\left(
\frac{\partial S}{\partial T}
\right)_H dT
+
\left(
\frac{\partial S}{\partial H}
\right)_T dH
\end{align*}

である。

ここで

\begin{align*}
C_H
=
T
\left(
\frac{\partial S}{\partial T}
\right)_H
\end{align*}

より、

\begin{align*}
\left(
\frac{\partial S}{\partial T}
\right)_H
=
\frac{C_H}{T}
\end{align*}

である。

また Maxwell 関係式を用いると、

\begin{align*}
\left(
\frac{\partial S}{\partial H}
\right)_T
=
\left(
\frac{\partial M}{\partial T}
\right)_H
\end{align*}

である。

したがって、

\begin{align*}
dS
=
\frac{C_H}{T}dT
+
\left(
\frac{\partial M}{\partial T}
\right)_H dH
\end{align*}

を得る。

以上より、示すべき式が導かれた。 【1-45ab9c】

\subsection{(ii)}

断熱変化では

\begin{align*}
dS=0
\end{align*}

である。

したがって前問の結果から、

\begin{align*}
0
=
\frac{C_H}{T}dT
+
\left(
\frac{\partial M}{\partial T}
\right)_H dH
\end{align*}

を得る。

よって、

\begin{align*}
\left(
\frac{\partial T}{\partial H}
\right)_S
=
-
\frac{T}{C_H}
\left(
\frac{\partial M}{\partial T}
\right)_H
\end{align*}

である。 \hfill (1)

ここで Curie の法則

\begin{align*}
M=\frac{CH}{T}
\end{align*}

を用いる。

温度で偏微分すると、

\begin{align*}
\left(
\frac{\partial M}{\partial T}
\right)_H
=
-\frac{CH}{T^2}
=
-\frac{M}{T}
\end{align*}

を得る。

これを式 (1) に代入すると、

\begin{align*}
\left(
\frac{\partial T}{\partial H}
\right)_S
=
\frac{M}{C_H}
\end{align*}

となる。

さらに問題5-2で求めた断熱条件

\begin{align*}
\frac{dT}{T}
=
\frac{H}{CC_M}dH
\end{align*}

を用いると、

\begin{align*}
\left(
\frac{\partial T}{\partial H}
\right)_S
=
\frac{TH}{CC_M}
\end{align*}

である。

また Curie の法則

\begin{align*}
M=\frac{CH}{T}
\end{align*}

より、

\begin{align*}
\frac{M}{C_H}
=
\frac{TH}{CC_M}
\end{align*}

が成り立つので、

\begin{align*}
\left(
\frac{\partial T}{\partial H}
\right)_S
=
\frac{T}{H}
\end{align*}

を得る。

したがって、

\begin{align*}
\left(
\frac{\partial T}{\partial H}
\right)_S
=
\frac{T}{H}
\end{align*}

である。 【1-45ab9c】

\subsection{(iii)}

前問の結果より、

\begin{align*}
\left(
\frac{\partial T}{\partial H}
\right)_S
=
\frac{T}{H}
>0
\end{align*}

である。

したがって断熱変化において

\begin{align*}
dH<0
\end{align*}

とすると、

\begin{align*}
dT
=
\left(
\frac{\partial T}{\partial H}
\right)_S
dH
<
0
\end{align*}

となる。

ゆえに外部磁場を弱くすると温度も低下する。

以上より、理想常磁性体に強い磁場をかけて磁化した後、断熱的に磁場を減少させると磁性体の温度は下降する。

これが断熱消磁冷却の原理である。 【1-45ab9c】
\section{6-8}

式

\begin{align*}
dU
=
TdS+HdM
\end{align*}

および Curie の法則

\begin{align*}
M=\frac{CH}{T}
\end{align*}

が成り立つ理想常磁性体を考える。磁場一定での熱容量を

\begin{align*}
C_H
=
T
\left(
\frac{\partial S}{\partial T}
\right)_H
\end{align*}

磁化一定での熱容量を

\begin{align*}
C_M
=
T
\left(
\frac{\partial S}{\partial T}
\right)_M
\end{align*}

とする。 【1-57d4ac】

\subsection{(i)}

まず

\begin{align*}
C_M
=
\left(
\frac{\partial U}{\partial T}
\right)_M
\end{align*}

である。

5-1 で示したように、理想常磁性体の内部エネルギーは温度のみの関数である：

\begin{align*}
U=U(T)。
\end{align*}

したがって、

\begin{align*}
C_M
=
\frac{dU(T)}{dT}
\end{align*}

となる。

よって、

\begin{align*}
C_M=C_M(T)
\end{align*}

であり、磁場や磁化には依存しない。

以上より、

\begin{align*}
C_M
\end{align*}

は温度のみの関数であることが示された。 【1-57d4ac】

\subsection{(ii)}

まず熱力学第一法則

\begin{align*}
dU
=
TdS+HdM
\end{align*}

を変形する：

\begin{align*}
TdS
=
dU-HdM。
\end{align*}

ここで、

\begin{align*}
dU
=
C_MdT
\end{align*}

であるから、

\begin{align*}
TdS
=
C_MdT-HdM
\end{align*}

を得る。

さらに Curie の法則

\begin{align*}
H=\frac{MT}{C}
\end{align*}

を代入すると、

\begin{align*}
TdS
=
C_MdT
-
\frac{MT}{C}dM。
\end{align*}

両辺を $T$ で割れば、

\begin{align*}
dS
=
\frac{C_M}{T}dT
-
\frac{M}{C}dM
\end{align*}

を得る。

以上より、

\begin{align*}
dS
=
\frac{C_M}{T}dT
-
\frac{M}{C}dM
\end{align*}

が示された。 【1-57d4ac】

\subsection{(iii)}

磁場一定熱容量

\begin{align*}
C_H
=
T
\left(
\frac{\partial S}{\partial T}
\right)_H
\end{align*}

を求める。

前問の結果

\begin{align*}
dS
=
\frac{C_M}{T}dT
-
\frac{M}{C}dM
\end{align*}

を用いる。

磁場一定では

\begin{align*}
dH=0
\end{align*}

であるから、Curie の法則

\begin{align*}
M=\frac{CH}{T}
\end{align*}

を微分して、

\begin{align*}
dM
=
-\frac{CH}{T^2}dT
=
-\frac{M}{T}dT
\end{align*}

を得る。

これを代入すると、

\begin{align*}
dS
&=
\frac{C_M}{T}dT
+
\frac{M^2}{CT}dT \\
&=
\frac{1}{T}
\left(
C_M+\frac{M^2}{C}
\right)dT。
\end{align*}

したがって、

\begin{align*}
\left(
\frac{\partial S}{\partial T}
\right)_H
=
\frac{1}{T}
\left(
C_M+\frac{M^2}{C}
\right)
\end{align*}

であり、

\begin{align*}
C_H
=
T
\left(
\frac{\partial S}{\partial T}
\right)_H
=
C_M+\frac{M^2}{C}
\end{align*}

を得る。

よって、

\begin{align*}
C_H-C_M
=
\frac{M^2}{C}
\end{align*}

が示された。 【1-57d4ac】

\subsection{(iv)}

前問で得た関係式

\begin{align*}
C_H
=
C_M+\frac{M^2}{C}
\end{align*}

を温度一定で磁場について偏微分する。

\begin{align*}
\left(
\frac{\partial C_H}{\partial H}
\right)_T
=
\frac{2M}{C}
\left(
\frac{\partial M}{\partial H}
\right)_T。
\end{align*}

Curie の法則より、

\begin{align*}
\left(
\frac{\partial M}{\partial H}
\right)_T
=
\frac{C}{T}
\end{align*}

である。

したがって、

\begin{align*}
\left(
\frac{\partial C_H}{\partial H}
\right)_T
=
\frac{2M}{T}。
\end{align*}

さらに Curie の法則

\begin{align*}
M=\frac{CH}{T}
\end{align*}

および

\begin{align*}
C_H-C_M
=
\frac{M^2}{C}
\end{align*}

を用いると、

\begin{align*}
C_H-C_M
=
\frac{CH^2}{T^2}。
\end{align*}

ここで両辺を $H$ で微分すると、

\begin{align*}
\left(
\frac{\partial C_H}{\partial H}
\right)_T
=
\frac{2CH}{T^2}。
\end{align*}

再び

\begin{align*}
M=\frac{CH}{T}
\end{align*}

を代入すれば、

\begin{align*}
\left(
\frac{\partial C_H}{\partial H}
\right)_T
=
\frac{2M}{T}
=
\frac{2CH}{T^2}
\end{align*}

を得る。

したがって、

\begin{align*}
\left(
\frac{\partial C_H}{\partial H}
\right)_T
=
\frac{2CH}{T^2}
\end{align*}

が成り立つ。

以上より、磁場一定熱容量は磁場とともに増加し、その磁場依存性は

\begin{align*}
\left(
\frac{\partial C_H}{\partial H}
\right)_T
=
\frac{2CH}{T^2}
\end{align*}

で与えられる。 【1-57d4ac】

\end{document}